% concatenated from iterated_graph_Laplacian.tex
\documentclass{article}
\usepackage{arxiv}
\usepackage[T1]{fontenc}
\usepackage[utf8]{inputenc}
\usepackage[english]{babel}
\usepackage{amsmath}
\usepackage{amsthm}
\usepackage{booktabs}
\usepackage{graphicx}
\usepackage{algorithm}
\usepackage{algorithmic}
\usepackage{hyperref}
\usepackage[capitalise,noabbrev]{cleveref}
\usepackage{sidecap}
\usepackage{amssymb}
\usepackage{mathrsfs}
\usepackage{mathtools}
\usepackage{accents}
\usepackage{dsfont}
\usepackage{diagbox,multirow}
\usepackage{hhline}
\usepackage{array}
\usepackage{xcolor}
\usepackage{colortbl}
\usepackage{enumerate}
\usepackage{bm}
\usepackage[shortlabels]{enumitem}
\usepackage{makecell}
\usepackage{float}
\usepackage{framed}
\usepackage{varwidth}
\usepackage[export]{adjustbox}
\usepackage{xparse}
\usepackage{csquotes}
\usepackage[nocompress]{cite}
\usepackage{subcaption}

\usepackage{etoolbox}
\apptocmd{\sloppy}{\hbadness 10000\relax}{}{}
\usepackage{amsbsy}

\makeatletter
\newcommand{\conditem}[2]{%
  \item[#2]%
  \protected@edef\@currentlabel{#2}%
  \label{#1}%
}
\makeatother
\newtheorem{theorem}{Theorem}[section]
\newtheorem{proposition}[theorem]{Proposition}
\newtheorem{lemma}[theorem]{Lemma}
\newtheorem{corollary}[theorem]{Corollary}
\newtheorem{definition}[theorem]{Definition}
\newtheorem{hypothesis}[theorem]{Hypothesis}
\theoremstyle{remark}
\newtheorem{remark}[theorem]{Remark}

\crefname{hypothesis}{Hypothesis}{Hypotheses}
\Crefname{hypothesis}{Hypothesis}{Hypotheses}
\Crefname{ALC@unique}{Line}{Lines}



\newcommand{\gt}{\bm{x}_{\mathrm{gt}}}
\newcommand{\N}{\mathbb{N}}
\newcommand{\R}{\mathbb{R}}

\newcommand{\opnormpq}[3]{\left\lVert #1 \right\rVert_{#2\to #3}}
\newcommand{\diag}{\operatorname{diag}}
\newcommand{\norm}[1]{\left\lVert #1 \right\rVert}
\newcommand{\ip}[2]{\left\langle #1,#2\right\rangle}
\newcommand{\Rreg}{\mathcal{R}}
\newcommand{\Ssol}{\mathcal{S}}
\newcommand{\Ncal}{\mathcal{N}}

\newcommand{\bx}{{\bm{x}}}
\newcommand{\by}{{\bm{y}}}
\newcommand{\bz}{{\bm{z}}}
\newcommand{\bw}{{\bm{w}}}
\newcommand{\bh}{{\bm{h}}}
\newcommand{\bu}{{\bm{u}}}
\newcommand{\bv}{{\bm{v}}}
\newcommand{\be}{{\bm{e}}}
\newcommand{\br}{{\bm{r}}}
\newcommand{\bs}{{\bm{s}}}
\newcommand{\bd}{{\bm{d}}}
\newcommand{\ba}{{\bm{a}}}
\newcommand{\bb}{{\bm{b}}}
\newcommand{\bg}{{\bm{g}}}
\newcommand{\bgt}{{\bm{x}_{\mathrm{gt}}}}
\newcommand{\bnoise}{{\bm{\eta}}}
\newcommand{\bone}{\mathbf{1}}

\title{Iterated graph Laplacian for image restoration problems}

\author{
Stefano Aleotti\\
Department of Science and High Technology (DiSAT)\\
University of Insubria, Como 22100, Italy\\
\texttt{stefano.aleotti@uninsubria.it}
\And
Davide Bianchi\\
School of Mathematics (Zhuhai)\\
Sun Yat-sen University, Zhuhai 518055, China\\
\texttt{bianchid@mail.sysu.edu.cn}
\And
Florian Bo\ss mann\\
Department of Mathematics\\
Harbin Institute of Technology, Harbin 150001, China\\
\texttt{f.bossmann@hit.edu.cn}
\And
Marco Donatelli\\
Department of Science and High Technology (DiSAT)\\
University of Insubria, Como 22100, Italy\\
\texttt{marco.donatelli@uninsubria.it}
\And
Pietro Maurino\\
Department of Science and High Technology (DiSAT)\\
University of Insubria, Como 22100, Italy\\
\texttt{pmaurino@uninsubria.it}
}

\date{}

\colorlet{siaminlinkcolor}{green!50!black}
\colorlet{siamexlinkcolor}{red!50!black}
\hypersetup{
  colorlinks=true,
  allcolors=siaminlinkcolor,
  urlcolor=siamexlinkcolor,
  pdftitle={Iterated graph Laplacian for Image Restoration Problems},
  pdfauthor={S. Aleotti, D. Bianchi, F. Bossmann, M. Donatelli, and P. Maurino}
}

\begin{document}

\maketitle

\begin{abstract}
We study the graph Laplacian operator as a regularizer in a generalized Tikhonov framework for linear ill-posed problems. The Laplacian is updated iteratively from the current reconstruction, so that progressively sharper structural information about the solution is fed into the regularization term. We introduce three schemes: a standard one that rebuilds the Laplacian from each new iterate; an error-equation scheme that, following the error-based formulation of iterated Tikhonov regularization, builds the Laplacian from an estimate of the reconstruction error rather than of the image itself; and a mixed scheme combining the two. We establish convergence of all three schemes for noisy data under a priori parameter and stopping rules as the noise level tends to zero. Numerical experiments in two-dimensional computed tomography and image deblurring show consistent gains in reconstruction quality and sharper recovery of fine details.
\end{abstract}

\paragraph{Funding.}
D. Bianchi is supported by the Startup Fund of Sun~Yat-sen~University.

\paragraph{Keywords.}
graph Laplacian regularization, image restoration, inverse problems, iterated Tikhonov, variational regularization

\paragraph{2020 Mathematics Subject Classification.}
65F10, 65F22, 65J20, 68U10, 94A08


\section{Introduction}
We are concerned with linear image restoration problems, whose model equation reads
\begin{equation*}\label{eq:Model_Equation}\tag{ME}
	K\bx = \by^\delta, 
\end{equation*}
where $\bx$ is the digital image to be recovered, $X$ is a finite dimensional space of images, and
$K \colon X \to Y$ is a discretized, inherently ill-posed linear operator modeling
a degradation such as blurring, tomographic projection, subsampling, or data masking, to name a few representative examples.

Given a fixed $\gt \in X$ and $\by\coloneqq K\gt$, we want to recover a good approximation of the ground truth $\gt$ from a noisy observation $\by^\delta$ of $\by$, that is,
\begin{equation*}
	\by^\delta \coloneqq \by + \bm{\eta}, \qquad \|\by^\delta - \by\|_Y\leq \delta,
\end{equation*}
where $\bm{\eta}$ is a random perturbation, typically unavoidable in practical scenarios, and $\delta>0$ is an estimate of the noise intensity.
On the one hand, the goal is to obtain a reliable approximation of the ground truth signal $\gt$ when only a noisy observation $\by^\delta$ is available. On the other hand, the unperturbed system \cref{eq:Model_Equation} with $\delta = 0$ may be underdetermined, the noisy data $\by^\delta$ may not belong to the range of $K$, or the (pseudo-)inverse of $K$ may be highly sensitive even to small noise levels $\delta > 0$. For these reasons, the inverse problem associated with \cref{eq:Model_Equation} is ill-posed and therefore requires regularization. For a comprehensive overview of ill-posed inverse problems and regularization techniques, we refer the reader to
\cite{engl1996regularization,hansen1998rank,scherzer2009variational}.

The ill-posedness of $K$ and the presence of noise force the introduction of a regularization strategy to solve our model equation. Among the various regularization methods in the literature, we consider here the $\ell^2$--$\ell^q$ regularization for~\cref{eq:Model_Equation} (see, e.g., \cite{CH14,LMRS15,HLMRS17}) that reads as
\begin{equation}\label{model_eq2}
	\bx^\delta_\alpha \in	\underset{\bx \in X}{\arg\min} \left\{ \frac{1}{2}\norm{K\bx - \by^\delta}_Y^2 + \frac{\alpha}{q}\norm{D \bx}_q^q \right\},
\end{equation} 
where $D$ is the regularization operator, a linear mapping satisfying $\ker(D)\cap\ker(K)=\{\boldsymbol{0}\}$; see \cite[Chapter 8]{engl1996regularization} and \cite{morozov2012methods}. In this context, the term $\frac{1}{2}\norm{K\bx-\by^\delta}_Y^2$ quantifies the fidelity of the reconstruction; the term $\norm{D\bx}_q^q$, with $q\in[1,2]$, serves as a regularizer; and $\alpha>0$ balances the trade-off between data fidelity and regularization.

Typical choices for the regularization operator $D$ include linear differential operators, see, e.g., \cite{ROF92,hansen2006deblurring,DR14}.
However, the last decade has seen an increasing interest in nonlocal models and techniques from graph theory \cite{gilboa2007nonlocal,gilboa2009nonlocal,Peyre2008nonlocal,arias2009variational}. This approach has been further investigated in the context of image deblurring and computed tomography (CT) problems in \cite{bianchi2021graph_approximation,buccini2021graph,bianchi2021graph,lou2010image,zhang2010bregmanized}.

In particular, in \cite{bianchi2025data}, the authors investigate a variational method for ill-posed problems, named graphLa\texttt{+}$\Psi$, which
embeds a graph Laplacian operator in the regularization term. The novelty of this method lies in constructing the graph Laplacian based on a preliminary approximation of the solution, which is
obtained using any existing reconstruction method $\Psi$ from the literature. The need for an initial reconstruction is due to the fact that a graph structure should incorporate features such as interfaces and discontinuities of the true solution. However, it was observed that generating the graph directly from the observed,
noisy data $\by^{\delta}$ yields poor results for imaging tasks such as
deblurring or computed tomography. This is unsurprising, since in general
$\by^{\delta}$ does not share the same features as $\gt$, and may not even lie
in the same domain.

This initial preprocessing step plays an important role in the quality of the
final reconstruction: the richer the detail encoded in the initial estimate,
the more effectively the graph Laplacian propagates that information throughout
the optimization. The role of this step has already been examined in depth
within the graphLa\texttt{+}$\Psi$ method, where several reconstruction
techniques, both standard and data-learned, were compared depending on the
underlying inverse problem. In the same spirit, and independently, iterative
schemes that update the graph weights along the iterations were proposed
in~\cite{Peyre2008nonlocal,zhang2010bregmanized,arias2009variational}.

What these approaches share is an iterative refinement principle: rather
than committing to a single initial estimate, one repeatedly improves it and
feeds the result back into the reconstruction. This principle is not specific to
graph-based regularization, and it has long been exploited for inverse problems
more broadly~\cite{nagy2026mixed}. A classical instance is the iterated Tikhonov
method~\cite[Chapter~6]{engl1996regularization}, which has been studied
extensively and extended to nonsmooth regularization terms as
well~\cite{huang2013nonstationary,aleotti2025nested,jin2012nonstationary,jin2014nonstationary,bianchi2023uniformly,bianchi2026iterated,bianchi2015iterated,bianchi2025convergence}.
 
 
In this work, we introduce three \emph{iterated graph Laplacian regularization
schemes}. In each of them, the reconstruction from which the graph Laplacian is
built is updated iteratively, so that the regularizer is driven by progressively
more accurate structural information about the solution.

The common building block is one graph Laplacian regularization step. Given a
reference image $\bz\in X$, we construct the associated graph Laplacian
$\Delta_{\bz}$ (see \Cref{sec:graph_theory}) and solve
\begin{equation}\label{eq:graphLa+}
\bx_{\alpha,\bz}^{\delta} \in
\underset{\bx \in X}{\arg\min}
\left\{
\tfrac{1}{2}\norm{K\bx - \by^\delta}_Y^2
+ \tfrac{\alpha}{q}\|\Delta_{\bz}\bx\|_q^q
\right\}.
\end{equation}
The three schemes differ in how this step is iterated.

The \emph{standard iterative scheme} feeds each new reconstruction back as the
reference for the next step: the graph Laplacian is recomputed from the current
iterate, \Cref{eq:graphLa+} is solved again, and the procedure is repeated. As
the iterates improve, so does the structural prior encoded by the graph
Laplacian.

The \emph{error-equation iterative scheme} takes a different viewpoint. Instead
of refining the solution directly, it refines the initial reconstruction $\bz$
by estimating the error. Following the iterated Tikhonov method, the
data-fidelity term $\tfrac12\norm{K\bx-\by^\delta}_Y^2$ is replaced by
$\tfrac12\norm{K\bh-\br^\delta}_Y^2$, where $\br^\delta=\by^\delta-K\bz$ is the
residual and $\bh$ the error correction to be computed. The graph Laplacian is now built from an approximation of the error $\bh$
rather than of the solution.

The \emph{mixed iterative scheme} combines the standard iterative approach with the
error-equation strategy: it begins with a small number of standard iterations to
recover the dominant structure, and then passes to the error-equation
formulation to sharpen the remaining fine-scale details. The motivation for this
combination lies in the complementary behavior of the two approaches.
Empirically, after a few iterations of the standard scheme the performance
metrics reach a plateau, so that further iterations yield no significant gain in
reconstruction quality. We interpret this as an indication that the graph
Laplacian built from the current iterate can no longer extract additional
structural information about the ground-truth solution $\gt$. Once this plateau
is reached, switching to the error-based formulation recovers further detail and
improves the reconstruction.

The standard iterative scheme was first proposed in~\cite{bianchi2024improved}, under the name \texttt{it-graphLa}$\Psi$,
for acoustic impedance inversion in seismic exploration, but only as a heuristic
refinement, without any convergence analysis. A related idea was developed
independently in~\cite{bajpai2026convergence}: there too the graph Laplacian is
recomputed from the current iterate, but the reconstruction is updated by an
explicit Landweber-type step with an extra graph term,
\[
\bx_{k+1}^{\delta}
= \bx_k^{\delta}
- \alpha_k^{\delta}\, K^*\!\left(K\bx_k^{\delta} - \by^\delta\right)
- \beta_k^{\delta}\, \Delta_{\bx_k^{\delta}}\bx_k^{\delta},
\]
and the iteration is terminated by the Morozov discrepancy principle. Avoiding a
variational subproblem at each update makes that scheme computationally lighter,
since one step costs essentially a forward/adjoint evaluation plus sparse graph
operations. Its convergence theory, however, is established in a fixed
finite-dimensional setting: the monotonicity constants and the admissible
step-size ranges depend on the discretization and on the radius of a ball,
centered at the initial reconstruction, which is assumed to contain the exact
solution.

Our analysis follows a different route. We retain the variational
formulation~\cref{eq:graphLa+} and, beyond the standard scheme, introduce and
analyze the error-equation and mixed schemes, which have no counterpart
in~\cite{bianchi2024improved,bajpai2026convergence}. We prove convergence of all
three schemes for noisy data, as the noise level tends to zero, under an a priori
parameter/stopping rule. The error-equation and mixed schemes
additionally require a uniform bound on the iterates, which holds automatically
under the grayscale box constraint $0\leq\bx\leq 1$. Throughout, all constants are
tracked explicitly and shown to be independent of the discretization level, that
is, of the number of pixels $N$, and of the reference image $\bz$.

In particular, these
results rest on minimal assumptions: the existence of a uniformly bounded exact
solution, and a joint coercivity condition coupling the forward operator and the
graph Laplacian that goes back to
Morozov~\cite{morozov2012methods} (see \Cref{sec:preliminary_estimates}). Both are required only to hold uniformly in
$N$ and $\bz$, and are otherwise standard. 

The paper is organized as follows. \Cref{sec:graph_theory} presents the discrete
image model, the graph construction, and the dimension-free graph estimates used
throughout the analysis. \Cref{sec:preliminary_estimates} introduces the two main assumptions and the 
fixed-reference variational problem, and proves the associated well-posedness and
stability results. \Cref{sec:iterschemes} defines the standard, error-equation,
and mixed iterative graph Laplacian schemes. \Cref{sec:convergence_standard}
analyzes the regularization properties of the standard scheme, and \Cref{sec:convergence_error}
treats the error-equation and mixed schemes. \Cref{sec:numerical_results} reports
extensive numerical experiments on image deblurring and computed tomography.
Finally, \Cref{sec:appendix-coercivity} verifies the joint coercivity assumption for the
forward operators used in the numerical experiments.




 
\section{Graph setting and graph estimates}\label{sec:graph_theory}

The variational approach considered in \Cref{eq:graphLa+} and all the iterative strategies introduced later rely on the graph Laplacian operator as a regularization tool. We therefore start by describing the discrete image model, the theoretical framework, and the graph construction.

Let $P$ be a finite set of pixels, with $N=|P|$, and let 
\[
X\coloneqq\R^{P}.
\]
We consider grayscale images as functions on the pixel set into $[0,1]$, that is
\[
\bx\colon P \to [0,1].
\]
In particular, the set of grayscale images is a proper subset of $X$. For $1\le q<\infty$ we use the normalized discrete norm
\[
\norm{\bx}_q\coloneqq\left(\frac1N\sum_{i\in P}|\bx(i)|^q\right)^{1/q},
\qquad \bx\in X,
\]
and
\[
\norm{\bx}_\infty\coloneqq\max_{i\in P}|\bx(i)|.
\]
The data space $Y$ is assumed to be a Hilbert space, with norm $\norm{\cdot}_Y$ induced by its inner product. In the applications that we consider, $Y$ is a finite-dimensional normalized Euclidean space, as for $X$. 

The model forward map is a linear operator between the Banach spaces just defined,
\begin{equation}\label{eq:model_operator}
K\colon (X, \|\cdot\|_q) \to (Y, \|\cdot\|_Y).
\end{equation}
We will often omit the space norms, since they will be clear from the context, and we write $\|A\|$ for the operator norm.

For a comprehensive introduction to the graph-based framework adopted here, we refer to \cite{keller2021graphs}. 

\begin{definition}[Graph and graph Laplacian]
A graph over $P$ is a symmetric function $\omega\colon P\times P\to[0,\infty)$ such that $\omega(i,i)=0$ for every $i\in P$. Each pair $(i,j)$ such that $\omega(i,j)>0$ is called an \emph{edge}.

The graph Laplacian $\Delta\colon X\to X$ associated with a graph $\omega$ is defined by the action
\begin{equation*}
(\Delta \bu)(i)\coloneqq\sum_{j\in P}\omega(i,j)\bigl(\bu(i)-\bu(j)\bigr),
\qquad i\in P.
\end{equation*}
\end{definition}

In the numerical experiments of \Cref{sec:numerical_results} and in the theory below, the graph is built from a reference image using a local Gaussian construction.

We first fix the spatial geometry. Throughout, we identify each pixel $i\in P$ with its position $(i_1,i_2)\in\mathbb{Z}^2$ on the two-dimensional pixel grid, so that for $i,j\in P$
\[
\operatorname{dist}(i,j)\coloneqq|i_1-j_1|+|i_2-j_2|
\]
denotes the Manhattan distance between the two pixel locations. This spatial distance is only used to select neighbors and is unrelated to the normalized norm $\|\cdot\|_1$ on $X$ introduced above. Different geometric distances may be used; however, for simplicity, we use only the Manhattan distance in this manuscript.

Fix now a spatial radius $R\in\N$ and a Gaussian parameter $\sigma>0$. For $i\in P$ define 
\[
\Ncal_R(i)\coloneqq\{j\in P \mid 0< \operatorname{dist}(i,j) \le R\}.
\]
Observe that the \emph{spatial degree} $|\Ncal_R(i)|$ is uniformly bounded:
\begin{equation}\label{eq:spatial-degree_uniform_norm}
d_R\coloneqq\sup_N\max_{i\in P}|\Ncal_R(i)|<\infty .
\end{equation}


Fix also a contrast cap $B_{\rm c}>0$, independent of $N$, and set
\begin{equation}\label{eq:contrast-cap}
d_{B_{\rm c}}(a,b)\coloneqq\min\{|a-b|,B_{\rm c}\},
\qquad a,b\in\R.
\end{equation}



Given a reference image $\bz\in X$, define
\begin{equation}\label{eq:gaussian-weights}
\omega_\bz(i,j)
\coloneqq\mathds{1}_{\{0<\operatorname{dist}(i,j)\le R\}}
\exp\!\left(-\frac{d_{B_{\rm c}}(\bz(i),\bz(j))^2}{\sigma^2}\right),
\qquad i,j\in P.
\end{equation}
The graph Laplacian induced by $\bz$ is
\begin{equation*}
(\Delta_\bz \bu)(i)
\coloneqq\sum_{j\in \Ncal_R(i)} \omega_\bz(i,j)\bigl(\bu(i)-\bu(j)\bigr),
\qquad \bu\in X .
\end{equation*}
We will often call $\bz$ the \emph{graph input} or the \emph{reference image}.


Let us remark that the contrast cap $B_{\rm c}$ acts only on the contrast range and imposes no
constraint on the graph input. It has no effect whenever
 the values of the graph input lie in an interval of length at most
$B_{\rm c}$; in particular, choosing $B_{\rm c}\ge1$ leaves ordinary
$[0,1]$-valued grayscale graph inputs unchanged. 

For $q\in[1,2]$, define
\begin{equation}\label{def:graph-regularizer}
\Rreg_\bz^q(\bu)\coloneqq\frac1q\norm{\Delta_\bz \bu}_q^q .
\end{equation}


\subsection{Dimension-free graph estimates}\label{ssec:graph-estimates}

The next estimates are consequences of the local Gaussian construction and the fixed radius $R$. 

\begin{lemma}\label{lem:graph-bounds}
For every $q\in[1,\infty]$, every $N$, every reference image $\bz\in X$, and every $\bu\in X$,
\begin{equation}\label{eq:Delta-uniform-bound}
\norm{\Delta_\bz \bu}_q\le 2d_R\norm{\bu}_q.
\end{equation}
Moreover, for all $\bx,\bz,\bu\in X$,
\begin{equation}\label{eq:Delta-Lipschitz}
\norm{(\Delta_\bx-\Delta_\bz)\bu}_q
\le L_\Delta\,\norm{\bx-\bz}_\infty\,\norm{\bu}_q,
\qquad
L_\Delta\coloneqq\frac{4d_R}{\sigma}\sqrt{\frac{2}{\mathrm{e}}} .
\end{equation}
The constants $2d_R$ and $L_\Delta$ are independent of $N$.
\end{lemma}

\begin{proof}
Let $A_\bz \colon X \to X$ be the adjacency operator
\[
(A_\bz \bu)(i)\coloneqq\sum_{j\in\Ncal_R(i)}\omega_\bz(i,j)\bu(j),
\]
and let $D_\bz \colon X \to X$ be the degree operator
\[
(D_\bz \bu)(i) \coloneqq \operatorname{deg}_\bz(i)\bu(i), \qquad \operatorname{deg}_\bz(i)\coloneqq\sum_{j\in\Ncal_R(i)}\omega_\bz(i,j).
\]
Since $0\le \omega_\bz(i,j)\le1$ and every row and column has at most $d_R$ nonzero entries by~\Cref{eq:spatial-degree_uniform_norm}, we have
\[
\norm{A_\bz}\le d_R,
\qquad
\norm{D_\bz}\le d_R
\]
for every $q\in[1,\infty]$. Since $\Delta_\bz=D_\bz-A_\bz$, \Cref{eq:Delta-uniform-bound} follows.

 

For the Lipschitz estimate, set
\[
g_{B_{\rm c}}(t)\coloneqq
\exp\!\left(-\frac{\min\{t,B_{\rm c}\}^2}{\sigma^2}\right),
\qquad t\ge0.
\]
The function $g_{B_{\rm c}}$ is globally Lipschitz and
\[
\operatorname{Lip}(g_{B_{\rm c}})
\le\frac1\sigma\sqrt{\frac{2}{\mathrm{e}}}.
\]
For every edge $(i,j)$,
\[
\begin{aligned}
|\omega_\bx(i,j)-\omega_\bz(i,j)|
&\le \frac1\sigma\sqrt{\frac{2}{\mathrm{e}}}
\bigl||\bx(i)-\bx(j)|-|\bz(i)-\bz(j)|\bigr|  \\
&\le \frac1\sigma\sqrt{\frac{2}{\mathrm{e}}}
\bigl(|\bx(i)-\bz(i)|+|\bx(j)-\bz(j)|\bigr) \\
&\le \frac{2}{\sigma}\sqrt{\frac{2}{\mathrm{e}}}\,\norm{\bx-\bz}_\infty.
\end{aligned}
\]
The difference $\Delta_\bx-\Delta_\bz$ is again a difference between a degree operator and an adjacency operator. Applying the same row-column bound as above gives \Cref{eq:Delta-Lipschitz}.
\end{proof}

\begin{corollary}\label{cor:R-bounds}
For every $q\in[1,2]$,
\begin{equation}\label{eq:R-uniform-bound}
0\le \Rreg_\bz^q(\bu)\le \frac{(2d_R)^q}{q}\norm{\bu}_q^q .
\end{equation}
Furthermore, there exists a constant
\[
C_{R,q}\coloneqq2(2d_R)^{q-1}L_\Delta
\]
independent of $N$ such that, for all $\bx,\bz,\bu\in X$,
\begin{equation}\label{eq:R-Lipschitz-guide}
\bigl|\Rreg_\bx^q(\bu)-\Rreg_\bz^q(\bu)\bigr|
\le C_{R,q}\,\norm{\bx-\bz}_\infty\,\norm{\bu}_q^q .
\end{equation}
Consequently, if $\bz_m\to \bz$ in $\norm{\cdot}_\infty$ and $\bu_m\to \bu$ in $X$, then
\begin{equation}\label{eq:R-continuity}
\Rreg_{\bz_m}^q(\bu_m)\to \Rreg_\bz^q(\bu).
\end{equation}
\end{corollary}

\begin{proof}
Estimate \Cref{eq:R-uniform-bound} follows immediately from \Cref{eq:Delta-uniform-bound}.

Let $\ba\coloneqq\Delta_\bx \bu$ and $\bb\coloneqq\Delta_\bz \bu$. For $q\in[1,2]$,
\[
\frac1q\bigl|\norm{\ba}_q^q-\norm{\bb}_q^q\bigr|
\le \bigl(\norm{\ba}_q^{q-1}+\norm{\bb}_q^{q-1}\bigr)
\norm{\ba-\bb}_q.
\]
 Using \Cref{lem:graph-bounds},
\[
\norm{\ba}_q^{q-1}+\norm{\bb}_q^{q-1}
\le 2(2d_R)^{q-1}\norm{\bu}_q^{q-1},
\qquad
\norm{\ba-\bb}_q\le L_\Delta\norm{\bx-\bz}_\infty\norm{\bu}_q,
\]
and \Cref{eq:R-Lipschitz-guide} follows.

For \Cref{eq:R-continuity}, write
\[
\begin{aligned}
|\Rreg_{\bz_m}^q(\bu_m)-\Rreg_\bz^q(\bu)|
&\le |\Rreg_{\bz_m}^q(\bu_m)-\Rreg_\bz^q(\bu_m)|
+|\Rreg_\bz^q(\bu_m)-\Rreg_\bz^q(\bu)|.
\end{aligned}
\]
The first term tends to zero by \Cref{eq:R-Lipschitz-guide}, since $(\|\bu_m\|_q)$ is bounded. The second term tends to zero by the continuity (uniform in $N$) of the fixed linear map $\Delta_\bz$ (see~\Cref{eq:Delta-uniform-bound}) and of the map $\bv\mapsto\norm{\bv}_q^q$.
\end{proof}

\section{Single iteration step and well-posedness}\label{sec:preliminary_estimates}
This section sets up the regularization framework and the main hypotheses before the iterative schemes are introduced. We then refine several stability estimates for the fixed-reference, or ``single-step'', variational problem involving the graph Laplacian, first established in~\cite{bianchi2025data}. These facts will be used repeatedly in the convergence analysis.

At discretization level $N$, given the model operator $K$ in~\Cref{eq:model_operator}, we recall the model equation~\cref{eq:Model_Equation}
\begin{equation*}
K\bx=\by^\delta,
\end{equation*}
where $\by^\delta$ is a noisy perturbation of the true data $\by$ such that
\begin{equation*}
\norm{\by^\delta-\by}_Y\le \delta.
\end{equation*}

\begin{hypothesis}[Existence of solutions]\label{hyp:consistency}
The observed data is the realization of a ground truth through the forward
operator, $\by=K\gt$, and the ground truth is uniformly bounded across
discretization levels: there is a constant $M_{\rm gt}>0$, independent of $N$, such that
\[
K\gt=\by,
\qquad
\norm{\gt}_2\le M_{\rm gt} .
\]
Since the normalized norms are nondecreasing in the exponent, it follows that
$\norm{\gt}_q\le M_{\rm gt}$ for every $q\in[1,2]$.
In particular, for every discretization level $N$, the solution set
\[
\Ssol(\by)\coloneqq\{\bx\in X:K\bx=\by\}
\]
is nonempty, since it contains $\gt$, and contains a uniformly bounded element.
For images in a fixed grayscale range, this bound holds automatically in the
normalized norms.
\end{hypothesis}

\begin{hypothesis}[Uniform joint coercivity]\label{hyp:joint-coercivity}
There exist constants $C_J>0$ and $M_K>0$, independent of $N$ and of the reference image $\bz$, such that
\begin{equation*}
\norm{K\bu}_Y\le M_K\norm{\bu}_2
\qquad\forall \bu\in X,
\end{equation*}
and
\begin{equation}\label{eq:joint-coercivity}
\norm{\bu}_2
\le C_J\Bigl(\norm{K\bu}_Y+\norm{\Delta_\bz \bu}_q\Bigr)
\qquad\forall \bu,\bz\in X .
\end{equation}
\end{hypothesis}

Condition~\cref{eq:joint-coercivity} is the graph Laplacian $\ell^q$ analogue of
the classical \emph{complementation condition} of Morozov~\cite{morozov2012methods}
(see also~\cite[Chapter~8]{engl1996regularization}). In the Hilbert-space case $q=2$, it is usually stated in the equivalent squared
form: there exists $\gamma>0$ such that
\[
\norm{K\bu}_Y^2+\norm{\Delta_\bz\bu}_2^2 \;\geq\; \gamma\,\norm{\bu}_2^2
\qquad\forall\,\bu\in X .
\]
The two formulations, for $q=2$, are equivalent, since $\tfrac{1}{\sqrt2}(a+b)\leq\sqrt{a^2+b^2}
\leq a+b$ for $a,b\geq0$: one passes from the squared form to
\Cref{eq:joint-coercivity} with $C_J=\gamma^{-1/2}$, and back with
$\gamma=(2C_J^2)^{-1}$.

A necessary requirement is that
the forward operator and the regularizer share no nontrivial kernel,
\[
\ker K\cap\ker\Delta_\bz=\{\boldsymbol 0\},
\]
and in the present finite-dimensional setting, this is also sufficient. What \Cref{hyp:joint-coercivity} demands
beyond this pointwise statement is that the bound holds \emph{uniformly} in the
discretization level $N$ and in the reference image $\bz$. This uniformity, rather
than the inequality at a fixed $N$, is what yields the dimension-free constants in
the convergence analysis.
The coercivity condition in \Cref{hyp:joint-coercivity} is then verified in \Cref{sec:appendix-coercivity} for some representative model operators $K$.



Now we introduce the definition of solutions for the model equation~\cref{eq:Model_Equation} with respect to the graph regularizer. It is standard in variational problems, see, for example, \cite[Definition~3.24]{scherzer2009variational}.

\begin{definition}[Regularizing solutions]\label{def:R-min-solution}
Fix $\bz\in X$ and define $\Rreg_\bz^q$ as per~\Cref{def:graph-regularizer}. A point $\bx^\star\in \Ssol(\by)$ is called an $\Rreg_\bz^q$-minimizing solution if
\begin{equation*}
\Rreg_\bz^q(\bx^\star)
=\min\{\Rreg_\bz^q(\bs):\bs\in\Ssol(\by)\}.
\end{equation*}
\end{definition}

The iterative schemes below (\Cref{sec:iterschemes}) update the reference image $\bz$ at every step, so the
limiting solution is not the minimizer for one fixed reference image but an
element of the collection of all minimizers obtained as the reference image
ranges over the exact solutions. We make this precise.

\begin{definition}[Graph-minimizing solutions]\label{def:Sgm}
The set of \emph{graph-minimizing solutions} of \Cref{eq:Model_Equation} is
\begin{equation}\label{eq:Sgm-def}
\mathcal S_{\rm gm}(\by)
\coloneqq
\bigcup_{\bz\in\Ssol(\by)}
\operatorname*{argmin}_{\bs\in\Ssol(\by)}\Rreg_\bz^q(\bs),
\end{equation}
that is, $\bu\in\mathcal S_{\rm gm}(\by)$ if and only if $\bu$ is an
$\Rreg_\bz^q$-minimizing solution in the sense of \Cref{def:R-min-solution} for
some reference $\bz\in\Ssol(\by)$.
\end{definition}

\begin{remark}\label{rem:Sgm}
   By construction $\mathcal S_{\rm gm}(\by)\subseteq\Ssol(\by)$. Under
\Cref{hyp:consistency,hyp:joint-coercivity} the set is
nonempty: for each fixed $\bz\in\Ssol(\by)$ the joint coercivity
\Cref{eq:joint-coercivity} makes the sublevel sets of $\Rreg_\bz^q$ within the
affine solution set $\Ssol(\by)$ bounded, so the inner $\operatorname*{argmin}$
is attained. The set reduces to a single point when $K$ is injective, since then
$\Ssol(\by)=\{\gt\}$. 
\end{remark}

\subsection{Single-step variational problem}\label{ssec:fixed-graph}
With the framework and hypotheses in place, we now study the single-step
problem~\cref{eq:graphLa+}, in which the reference image $\bz$ is held fixed.
This is the building block of all three iterative schemes. Each iteration updates
the reference $\bz$ (through the recomputed graph Laplacian), and the
error-equation and mixed schemes update the data $\by$ as well (through the
residual). We therefore need the single step to be well defined and to depend stably on both
$\by$ and $\bz$, with constants independent of the discretization level $N$.

The next \Cref{prop:existence} establishes that the fixed-reference functional is convex
and coercive, with sublevel sets bounded uniformly in $N$ and $\bz$; a minimizer
always exists, and is unique when $q>1$. \Cref{cor:one-step-stability} then shows
that this minimizer varies continuously with the data $\by$ and the reference
$\bz$, again uniformly in $N$. Both facts are used repeatedly in the convergence
analysis in \Cref{sec:convergence_standard,sec:convergence_error}.

For $q\in[1,2]$, $\alpha>0$, a reference image $\bz\in X$, and data $\by\in Y$,
define
\begin{subequations}\label{eq:fixed-graph-functional}
\begin{align}
\Phi_{\alpha,\bz}^q(\bx;\by)
&\coloneqq\frac12\norm{K\bx-\by}_Y^2+\alpha\Rreg_\bz^q(\bx),
\qquad \bx\in X ,
\label{eq:fixed-graph-functional-a}\\[6pt]
T^q_{\alpha,\bz}(\by)
&\coloneqq\operatorname*{argmin}_{\bx\in X}\;\Phi_{\alpha,\bz}^q(\bx;\by).
\label{eq:fixed-graph-functional-b}
\end{align}
\end{subequations}
This notation will be used with either exact data $\by$ or noisy data
$\by^\delta$.

\begin{proposition}\label{prop:existence}
Assume \Cref{hyp:joint-coercivity}. For every $N$, $\bz\in X$,
$\by\in Y$, $\alpha>0$, and $q\in[1,2]$, the functional
$\Phi_{\alpha,\bz}^q(\cdot;\by)$ is convex and coercive. More precisely, if
\[
\Phi_{\alpha,\bz}^q(\bx;\by)\le M,
\]
then
\begin{equation}\label{eq:sublevel-bound}
\norm{\bx}_2
\le C_J\left((2M)^{1/2}+\norm{\by}_Y
+\left(\frac{qM}{\alpha}\right)^{1/q}\right).
\end{equation}
Consequently, $\Phi_{\alpha,\bz}^q(\cdot;\by)$ admits at least one minimizer.
If $q>1$, this minimizer is unique.
\end{proposition}

\begin{proof}
Convexity follows from the convexity of the data term $\bx \mapsto \frac12\norm{K\bx - \by}_Y^2$ and of the regularization term
$\bx\mapsto\norm{\Delta_\bz \bx}_q^q$. Assume $\Phi_{\alpha,\bz}^q(\bx;\by)\le M$. Then
\[
\norm{K\bx-\by}_Y\le (2M)^{1/2},
\qquad
\norm{\Delta_\bz \bx}_q\le \left(\frac{qM}{\alpha}\right)^{1/q}.
\]
Hence $\norm{K\bx}_Y\le (2M)^{1/2}+\norm{\by}_Y$. Applying
\Cref{eq:joint-coercivity} gives \Cref{eq:sublevel-bound}. Thus the sublevel
sets are bounded with a bound independent of $N$ and $\bz$. Since $X$ is
finite-dimensional and $\Phi_{\alpha,\bz}^q(\cdot;\by)$ is continuous, a
minimizer exists; see, for example, \cite[Proposition 11.15]{bauschke2017convex}.

For uniqueness when $q>1$, let $\bx_1\ne\bx_2$ and set
$\bu\coloneqq\bx_1-\bx_2$. By \Cref{eq:joint-coercivity}, $K\bu$ and
$\Delta_\bz\bu$ cannot both vanish. If $K\bu\ne0$, the data term is strictly
convex on the segment joining $\bx_1$ and $\bx_2$. If $\Delta_\bz\bu\ne0$, then
$\norm{\Delta_\bz \cdot}_q^q$ is strictly convex for $q>1$, so the regularization term is
strictly convex on that segment. Therefore
$\Phi_{\alpha,\bz}^q(\cdot;\by)$ is strictly convex along every nontrivial
segment. Uniqueness then follows.
\end{proof}


\begin{corollary}
\label{cor:one-step-stability}
Assume \Cref{hyp:joint-coercivity}, and fix $q\in[1,2]$ and $\alpha>0$,
with $\Phi^q_{\alpha,\bz}$ and $T^q_{\alpha,\bz}$ as in~\Cref{eq:fixed-graph-functional}. Let
\[
\by_m\to\by \;\text{ in } Y,
\qquad
\norm{\bz_m-\bz}_\infty\to 0,
\]
and choose any
\[
\bx_m\in T^q_{\alpha,\bz_m}(\by_m).
\]
Then the following hold.

\begin{enumerate}[(i)]
\item The family $(\bx_m)$ is bounded uniformly in $N$, every cluster point
of $(\bx_m)$ is a minimizer of $\Phi_{\alpha,\bz}^q(\cdot;\by)$, and consequently
$(\bx_m)$ has a subsequence converging in $X$ to such a minimizer.
\item If $q>1$, the minimizer is unique, so the operator in
\Cref{eq:fixed-graph-functional-b} is single-valued, $\bx_m=T^q_{\alpha,\bz_m}(\by_m)$,
and the whole sequence converges,
\[
\bx_m\to T^q_{\alpha,\bz}(\by)\quad\text{in }X.
\]
\end{enumerate}
\end{corollary}
\begin{proof}
By minimality, testing against $\bx=0$ in~\Cref{eq:fixed-graph-functional-a} gives
\[
\Phi_{\alpha,\bz_m}^q(\bx_m;\by_m)
\le \Phi_{\alpha,\bz_m}^q(0;\by_m)
=\tfrac12\norm{\by_m}_Y^2 .
\]
Since $(\by_m)$ converges it is bounded, so the right-hand side is bounded
uniformly in $m$, and \Cref{prop:existence} yields a bound on
$(\bx_m)$ independent of both $m$ and $N$. As $X$ is finite-dimensional,
$(\bx_m)$ therefore admits cluster points. Let $\bx_{m_k}\to\bar\bx$ in $X$.
For every fixed $\bx\in X$, minimality of $\bx_{m_k}$ gives
\[
\Phi_{\alpha,\bz_{m_k}}^q(\bx_{m_k};\by_{m_k})
\le \Phi_{\alpha,\bz_{m_k}}^q(\bx;\by_{m_k}).
\]
The data-fidelity term in~\Cref{eq:fixed-graph-functional-a} is continuous in
$(\bx,\by)$, and by \Cref{cor:R-bounds} the map
$\bz\mapsto\Rreg_\bz^q(\bu)$ is continuous; passing to the limit along $k$ on both
sides (using $\by_{m_k}\to\by$ and $\bz_{m_k}\to\bz$) yields
\[
\Phi_{\alpha,\bz}^q(\bar\bx;\by)
\le \Phi_{\alpha,\bz}^q(\bx;\by)
\qquad\forall\,\bx\in X .
\]
Hence $\bar\bx$ minimizes $\Phi_{\alpha,\bz}^q(\cdot;\by)$, which proves~(i).

If $q>1$, \Cref{prop:existence} shows each functional is strictly
convex, hence has a unique minimizer; the $\operatorname*{argmin}$
in~\Cref{eq:fixed-graph-functional-b} is then single-valued, so
$\bx_m=T^q_{\alpha,\bz_m}(\by_m)$ and the limit problem has the unique minimizer
$T^q_{\alpha,\bz}(\by)$. By~(i), every subsequence of $(\bx_m)$ has a further
subsequence converging to this same limit, so by standard topological arguments the whole sequence converges,
which proves~(ii).
\end{proof}


\section{Iterated graph Laplacian regularization schemes}\label{sec:iterschemes}
In this section, we formally introduce the new iterative schemes based on the graph Laplacian defined in \Cref{sec:graph_theory}. These algorithms include the standard graph-based update, the error-equation strategy, and the mixed iterative method combining the previous two.

Since $T^q_{\alpha,\bz}$ is set-valued when minimizers are not unique, each time we write $\bu\in T^q_{\alpha,\bz}(\by)$ we mean that one admissible minimizer is selected.

\subsection{Standard iterative scheme}
Given a preliminary reconstruction $\bx_0^\delta \in X$ and a positive sequence $\{\alpha_n\}_{n\ge1}$, we define the iterative scheme
\begin{equation}\label{eq:standard-iter}
\bx_n^\delta \in T^q_{\alpha_n,\bx_{n-1}^\delta}(\by^\delta),
\qquad n \geq 1.
\end{equation}
Equivalently, at each iteration $\bx_n^\delta$ minimizes
\[
\frac12\norm{K\bx-\by^\delta}_Y^2+\alpha_n\Rreg_{\bx_{n-1}^\delta}^q(\bx)
\qquad\text{over }\bx\in X.
\]
Thus, the graph Laplacian operator is recomputed from the previous iterate $\bx_{n-1}^\delta$, and the regularization functional varies with $n$. The proposed strategy is described in \Cref{alg:standard-iter}.

\begin{algorithm}[h]
\caption{Standard iterative scheme}\label{alg:standard-iter}
\begin{algorithmic}[1]
\REQUIRE $K$, $\by^\delta$, $\bx_0^\delta$, $\{\alpha_n\}_{n\ge1}$
\FOR{$n = 1,2,\ldots$}
\STATE $\bx_n^\delta \in T^q_{\alpha_n,\bx_{n-1}^\delta}(\by^\delta)$
\ENDFOR
\end{algorithmic}
\end{algorithm}

The rationale behind this strategy is as follows. The graph Laplacian can effectively encode prior information about the ground truth solution provided that the reconstruction used to build it is a sufficiently accurate approximation of $\gt$. Even when the initial reconstruction $\bx_0^\delta$ is a poor approximation, the first iterate $\bx_1^\delta$ is expected to improve upon it. Consequently, constructing the graph Laplacian using $\bx_1^\delta$ yields a linear operator that is better suited as a regularizer. As the iterations proceed, the graph Laplacian progressively incorporates increasingly accurate structural information about the solution, leading to more effective regularization and, ultimately, to an improved final reconstruction.

\subsection{Error-equation iterative scheme}\label{sec:erreq}
A second approach is based on iterative refinement through the error equation. Starting from an initial reconstruction $\bx_0^\delta$, one may write a correction $\bh=\bx-\bx_0^\delta$. Then
\[
K\bh = \by^\delta-K\bx_0^\delta\eqqcolon \br_0^\delta
\]
is the residual equation for the first correction. In variational form, this leads to
\[
\underset{\bh\in X}{\operatorname{argmin}}
\left\{
\tfrac12\norm{K\bh-\br_0^\delta}_Y^2
+\alpha_1\Rreg_{\bh_0^\delta}^q(\bh)
\right\},
\]
where $\bh_0^\delta\in X$ is a prescribed image-space approximation of the first error, used only to construct the graph Laplacian in the first error step.

Once the data fidelity term is expressed through the error equation $K\bh=\br$, the graph Laplacian should be constructed from an image carrying information about the expected error rather than from an approximation of the true solution itself. The error-equation iteration therefore updates both the reconstruction and the error graph input. With
\[
\br_n^\delta\coloneqq\by^\delta-K\bx_n^\delta,
\]
the scheme is
\[
\left\{
\begin{aligned}
\bh_{n+1}^\delta
&\in T^q_{\alpha_{n+1},\bh_n^\delta}(\br_n^\delta),\\
\bx_{n+1}^\delta
&=\bx_n^\delta+\bh_{n+1}^\delta,\\
\br_{n+1}^\delta
&=\br_n^\delta-K\bh_{n+1}^\delta,
\end{aligned}
\right.
\qquad n\ge0.
\]
The algorithm is summarized in \Cref{alg:error-iter}.

\begin{algorithm}[h!]
\caption{Error-equation iterative scheme}\label{alg:error-iter}
\begin{algorithmic}[1]
\REQUIRE $K$, $\by^\delta$, $\bx_0^\delta$, $\bh_0^\delta$, $\{\alpha_n\}_{n\ge1}$
\STATE $\br_0^\delta=\by^\delta-K\bx_0^\delta$
\FOR{$n = 0,1,2,\ldots$}
\STATE $\bh_{n+1}^\delta \in T^q_{\alpha_{n+1},\bh_n^\delta}(\br_n^\delta)$
\STATE $\bx_{n+1}^\delta=\bx_n^\delta+\bh_{n+1}^\delta$
\STATE $\br_{n+1}^\delta=\br_n^\delta-K\bh_{n+1}^\delta$
\ENDFOR
\end{algorithmic}
\end{algorithm}

\subsection{Mixed iterative scheme}\label{sec:mix}
The mixed iterative scheme combines the standard approach in~\Cref{alg:standard-iter} with the error-equation strategy in~\Cref{alg:error-iter}. The motivation for combining these two approaches lies in their complementary behavior. Empirically, after a few iterations of the standard scheme, we observe that the performance metrics reach a plateau, meaning that further iterations do not lead to any significant improvement in the reconstruction quality. This phenomenon can be interpreted as an indication that the graph Laplacian, constructed from the current iterate, is no longer able to extract additional structural information about the ground truth solution $\gt$. Once this plateau is reached, switching to the error-based formulation can further improve the reconstruction by recovering additional details. 


\begin{algorithm}[t]
\caption{Mixed iterative scheme}\label{alg:mixed-iter}
\begin{algorithmic}[1]
\REQUIRE $K$, $\by^\delta$, $\bx_0^\delta$, $\{\alpha_n\}_{n\ge1}$, $n_{\mathrm{std}}$
\STATE \textbf{Standard phase}
\FOR{$n = 1$ \TO $n_{\mathrm{std}}$}
\STATE $\bx_n^\delta \in T^q_{\alpha_n,\bx_{n-1}^\delta}(\by^\delta)$
\ENDFOR
\STATE $\br_{n_{\mathrm{std}}}^\delta=\by^\delta-K\bx_{n_{\mathrm{std}}}^\delta$
\STATE Choose an image-space initialization $\bg_0^\delta$
\STATE \textbf{Error phase}
\FOR{$k = 0,1,2,\ldots$}
\STATE $\bg_{k+1}^\delta\in T^q_{\alpha_{n_{\mathrm{std}}+k+1},\bg_k^\delta}(\br_{n_{\mathrm{std}}+k}^\delta)$
\STATE $\bx_{n_{\mathrm{std}}+k+1}^\delta=\bx_{n_{\mathrm{std}}+k}^\delta+\bg_{k+1}^\delta$
\STATE $\br_{n_{\mathrm{std}}+k+1}^\delta=\br_{n_{\mathrm{std}}+k}^\delta-K\bg_{k+1}^\delta$
\ENDFOR
\end{algorithmic}
\end{algorithm}


More precisely, fix the number of standard iterations
$n_{\mathrm{std}}\in\mathbb{N}_{>0}$. We first run $n_{\mathrm{std}}$ steps
of~\Cref{alg:standard-iter}, producing $\bx_{n_{\mathrm{std}}}^\delta$, and then
continue with the error-equation iteration of~\Cref{alg:error-iter} initialized at
this point. We set the residual
\[
\br_{n_{\mathrm{std}}}^\delta\coloneqq\by^\delta-K\bx_{n_{\mathrm{std}}}^\delta ,
\]
and choose an initial error-stage graph input $\bg_0^\delta\in X$.
Since this choice is problem-dependent, in the numerical experiments we use the
following strategies:
\begin{itemize}
    \item \emph{Deblurring problem.} Since the residual belongs to the same image
    space as the reconstruction, we choose
    \begin{equation}\label{eq:g0res}
    \bg_0^\delta=\br_{n_{\mathrm{std}}}^\delta
    =\by^\delta-K\bx_{n_{\mathrm{std}}}^\delta .
    \end{equation}
    \item \emph{CT problem.} We apply filtered backprojection (FBP) to the
    residual $\br_{n_{\mathrm{std}}}^\delta$ to obtain an image-space approximation
    of the error.
    \item \emph{General approach.} A problem-independent alternative computes the
    initialization via Tikhonov regularization,
    \[
    \bg_0^\delta \in
    \underset{\bw \in X}{\operatorname{argmin}}
    \left\{
    \tfrac{1}{2}\norm{K\bw-\br_{n_{\mathrm{std}}}^\delta}_Y^2
    + \lambda \norm{\bw}_2^2
    \right\},
    \]
    where $\lambda>0$ can be selected by an automatic rule such as generalized
    cross-validation (GCV).
\end{itemize}
Finally, we perform a prescribed number of error-equation steps as
in~\Cref{alg:error-iter}, building the graph Laplacian from
$\bg_0^\delta$ at the first error step. Usually a single iteration
suffices to recover a few additional sharp edges. The overall procedure is summarized in~\Cref{alg:mixed-iter}.




\section{Regularization properties of \texorpdfstring{\Cref{alg:standard-iter}}{the standard iterative scheme}}\label{sec:convergence_standard}
In this section, we present the theoretical results concerning the standard iterative scheme.

\subsection{Finite-iteration stability}
We now show that the one-step stability of \Cref{cor:one-step-stability}
propagates, by induction, through any fixed number of iterations of the standard
scheme: if the data and the initial reconstruction
converge, so do the first $J$ iterates, with bounds independent of the
discretization level $N$. 
\begin{theorem}\label{thm:finite-stability}
Assume \Cref{hyp:joint-coercivity}. Fix $J\in\N$ and positive parameters $\alpha_1,\ldots,\alpha_J$. Let
\[
\by_m\to \by \quad\text{in }Y,
\qquad
\bx_{0,m}\to \bx_0 \quad\text{in }\norm{\cdot}_\infty.
\]
For each $m$, let $\bx_{n,m}$ be generated by \Cref{eq:standard-iter} with data $\by_m$ and initial point $\bx_{0,m}$, for $n=1,\ldots,J$.

Then the family of all such iterates is bounded, with bounds independent of $N$. Moreover, every subsequence indexed by $(m_k)$ has a further subsequence indexed by $(m_{k_j})$ such that, for each $n=1,\ldots,J$,
\[
\bx_{n,m_{k_j}}\to \bar \bx_n,
\]
where $(\bar \bx_n)_{n=1}^J$ is a sequence generated by the exact limiting data $\by$ and initial point $\bx_0$, namely
\[
\bar \bx_n\in\operatorname*{argmin}_{\bx\in X}
\left\{
\frac12\norm{K\bx-\by}_Y^2
+\alpha_n\Rreg_{\bar \bx_{n-1}}^q(\bx)
\right\},
\qquad \bar \bx_0=\bx_0.
\]
If the minimizer is unique at every limiting step, in particular if $q>1$, then the whole sequence converges:
\[
\bx_{n,m}\to \bar \bx_n,
\qquad n=1,\ldots,J.
\]
\end{theorem}

\begin{proof}
We argue by induction on $n$. For $n=1$, the claim is exactly \Cref{cor:one-step-stability}, applied with data $\by_m\to\by$, reference images $\bz_m=\bx_{0,m}$, and limit reference $\bz=\bx_0$.

Now assume the claim holds for $n-1$, so along any convergent subsequence $\bx_{n-1,m_k}\to\bar\bx_{n-1}$. Apply \Cref{cor:one-step-stability} once more, this time with reference images $\bz_{m_k}=\bx_{n-1,m_k}$ and limit reference $\bz=\bar\bx_{n-1}$. It follows that the sequence $\bx_{n,m_k}$ is bounded, and every cluster point of it solves the $n$-th problem in the limit.

So along every subsequence, the cluster points are minimizers. This proves the stability statement. Moreover, if the limiting problem has a unique minimizer at each step, there is only one possible cluster point, so the whole sequence converges.
\end{proof}

\subsection{Convergence for noisy data}

Classical regularization theory identifies the noise-free limit of the
regularized solutions with a single minimum-regularization solution, and
whole-sequence convergence is recovered only under a uniqueness assumption on
that solution. Here the situation is different. The regularizer
$\Rreg_\bz^q$ carries a reference point $\bz$ that is itself produced by the
iteration. Therefore, distinct
subsequences may converge to distinct reference images $\bar\bx$, and each reference image
can select a different minimizer. The natural limit object is thus not a
point but the set $\mathcal S_{\rm gm}(\by)$ of graph-minimizing solutions
introduced in \Cref{def:Sgm}.

Accordingly, our main convergence statement is set-valued: the iterates stopped at $n=n(\delta)$
approach this set,
\[
\operatorname{dist}_2\bigl(\bx_{n(\delta)}^\delta,\mathcal S_{\rm gm}(\by)\bigr)
\to0
\qquad\text{as }\delta\downarrow0,
\]
where
\[
\operatorname{dist}_2(\bx,U)
\coloneqq
\inf_{\bu\in U}\norm{\bx-\bu}_2,
\]
with no uniqueness assumption and no source condition. It reduces to the usual
single-point convergence exactly when $\mathcal S_{\rm gm}(\by)$ is a singleton,
for instance when $K$ is injective. 

Before proceeding further with the convergence result, we need a couple of preliminary estimates.

\medskip
\noindent\textit{A priori estimates.}
Assume \Cref{hyp:consistency,hyp:joint-coercivity}, and let
$\norm{\by^\delta-\by}_Y\le\delta$. Let $(\bx_n^\delta)$ be generated by
\Cref{eq:standard-iter}. With $\gt\in\Ssol(\by)$ the bounded ground truth from
\Cref{hyp:consistency}, the inequality~\cref{eq:R-uniform-bound} gives
\begin{equation*}
\Rreg_{\bz}^q(\gt)\le C_0,
\qquad
C_0\coloneqq\frac{(2d_R)^q}{q}M_{\rm gt}^q ,
\qquad \bz\in X .
\end{equation*}
Therefore, by minimality of $\bx_n^\delta$, for every $n\ge1$,
\begin{equation*}
\begin{aligned}
\frac12\norm{K\bx_n^\delta-\by^\delta}_Y^2
+\alpha_n\Rreg_{\bx_{n-1}^\delta}^q(\bx_n^\delta)
&\le
\frac12\norm{K\gt-\by^\delta}_Y^2
+\alpha_n\Rreg_{\bx_{n-1}^\delta}^q(\gt)\\
&\le \frac12\delta^2+\alpha_n C_0 .
\end{aligned}
\end{equation*}
Consequently,
\begin{equation}\label{eq:noisy-res-bound}
\norm{K\bx_n^\delta-\by^\delta}_Y^2
\le \delta^2+2\alpha_n C_0,
\end{equation}
and
\begin{equation}\label{eq:noisy-reg-bound}
\Rreg_{\bx_{n-1}^\delta}^q(\bx_n^\delta)
\le C_0+\frac{\delta^2}{2\alpha_n}.
\end{equation}

\begin{theorem}
\label{thm:standard-vanishing-noise}
Assume \Cref{hyp:consistency,hyp:joint-coercivity}.
Let $(\alpha_n)$ be nonincreasing, $\alpha_n\downarrow0$, and let the stopping
index $n(\delta)$ be chosen so that
\begin{equation}\label{eq:apriori-rule}
\alpha_{n(\delta)}\to0,
\qquad
\frac{\delta^2}{\alpha_{n(\delta)}}\to0
\qquad\text{as }\delta\downarrow0.
\end{equation}
Let $\norm{\by^\delta-\by}_Y\le\delta$, and let $(\bx_n^\delta)$ be generated by
\Cref{eq:standard-iter} with data $\by^\delta$. Then
\begin{equation}\label{eq:apriori-residual-convergence}
\norm{K\bx_{n(\delta)}^\delta-\by}_Y\to0
\qquad\text{as }\delta\downarrow0.
\end{equation}
Moreover, the family $(\bx_{n(\delta)}^\delta)$ is bounded in $X$, with
bounds independent of $N$, and
\begin{equation}\label{eq:set-convergence-Sgm}
\operatorname{dist}_2
\bigl(\bx_{n(\delta)}^\delta,\mathcal S_{\rm gm}(\by)\bigr)
\to0
\qquad\text{as }\delta\downarrow0.
\end{equation}
In particular,
\begin{equation}\label{eq:set-convergence-S}
\operatorname{dist}_2
\bigl(\bx_{n(\delta)}^\delta,\Ssol(\by)\bigr)
\to0
\qquad\text{as }\delta\downarrow0.
\end{equation}
\end{theorem}

\begin{proof}
By \Cref{rem:Sgm} the set $\mathcal S_{\rm gm}(\by)$ is nonempty, so the distance
in \Cref{eq:set-convergence-Sgm} is well defined.

Set
\[
m_\delta\coloneqq n(\delta).
\]
From \Cref{eq:apriori-rule} and $\alpha_n\downarrow0$ we have
$m_\delta\to\infty$ as $\delta\downarrow0$. Moreover,
\[
\alpha_{m_\delta-1}\to0,
\qquad
\frac{\delta^2}{\alpha_{m_\delta-1}}
\le
\frac{\delta^2}{\alpha_{m_\delta}}
\to0 ,
\]
where the inequality uses the monotonicity of $(\alpha_n)$.

By \Cref{eq:noisy-res-bound}, applied with $n=m_\delta$,
\[
\norm{K\bx_{m_\delta}^\delta-\by}_Y
\le
\norm{K\bx_{m_\delta}^\delta-\by^\delta}_Y
+\norm{\by^\delta-\by}_Y
\le
\bigl(\delta^2+2\alpha_{m_\delta}C_0\bigr)^{1/2}
+\delta .
\]
Therefore
\[
\norm{K\bx_{m_\delta}^\delta-\by}_Y\to0 ,
\]
which proves \Cref{eq:apriori-residual-convergence}. The same argument with
$n=m_\delta-1$ gives
\begin{equation}\label{eq:predecessor-residual-convergence}
\norm{K\bx_{m_\delta-1}^\delta-\by}_Y\to0 .
\end{equation}

Next, \Cref{eq:noisy-reg-bound}, used with $n=m_\delta$ and $n=m_\delta-1$,
gives uniform bounds on
\[
\Rreg_{\bx_{m_\delta-1}^\delta}^q(\bx_{m_\delta}^\delta),
\qquad
\Rreg_{\bx_{m_\delta-2}^\delta}^q(\bx_{m_\delta-1}^\delta).
\]
Together with \Cref{eq:apriori-residual-convergence},
\Cref{eq:predecessor-residual-convergence}, and the joint coercivity condition
\Cref{eq:joint-coercivity}, this gives uniform bounds on
$\bx_{m_\delta}^\delta$ and $\bx_{m_\delta-1}^\delta$, with constants
independent of $N$.

We now prove \Cref{eq:set-convergence-Sgm}. Suppose, by contradiction, that it
fails. Then there exist $\varepsilon>0$ and a sequence $\delta_j\downarrow0$
such that, with $m_j\coloneqq n(\delta_j)$,
\begin{equation}\label{eq:dist-contradiction}
\operatorname{dist}_2
\bigl(\bx_{m_j}^{\delta_j},\mathcal S_{\rm gm}(\by)\bigr)
\ge\varepsilon
\qquad\forall j .
\end{equation}
By the uniform bounds just proved, finite-dimensional compactness gives, after
passing to a subsequence,
\[
\bx_{m_j}^{\delta_j}\to\bx^\star,
\qquad
\bx_{m_j-1}^{\delta_j}\to\bar\bx .
\]
The residual convergence gives
\[
\bx^\star,\bar\bx\in\Ssol(\by).
\]

It remains to identify $\bx^\star$. Let $\bs\in\Ssol(\by)$ be arbitrary. By
minimality of $\bx_{m_j}^{\delta_j}$ and $K\bs=\by$,
\[
\Rreg_{\bx_{m_j-1}^{\delta_j}}^q(\bx_{m_j}^{\delta_j})
\le
\Rreg_{\bx_{m_j-1}^{\delta_j}}^q(\bs)
+\frac{\delta_j^2}{2\alpha_{m_j}} .
\]
Passing to the limit, using \Cref{cor:R-bounds} and
\Cref{eq:apriori-rule}, yields
\[
\Rreg_{\bar\bx}^q(\bx^\star)
\le
\Rreg_{\bar\bx}^q(\bs)
\qquad
\forall \bs\in\Ssol(\by).
\]
Thus
\[
\bx^\star\in
\operatorname*{argmin}_{\bs\in\Ssol(\by)}
\Rreg_{\bar\bx}^q(\bs).
\]
Since $\bar\bx\in\Ssol(\by)$, the definition \cref{eq:Sgm-def} gives
$\bx^\star\in\mathcal S_{\rm gm}(\by)$.

But then
\[
\operatorname{dist}_2
\bigl(\bx_{m_j}^{\delta_j},\mathcal S_{\rm gm}(\by)\bigr)
\le
\norm{\bx_{m_j}^{\delta_j}-\bx^\star}_2
\to0,
\]
contradicting \Cref{eq:dist-contradiction}. Hence
\Cref{eq:set-convergence-Sgm} holds. Since
$\mathcal S_{\rm gm}(\by)\subseteq\Ssol(\by)$,
\Cref{eq:set-convergence-S} follows immediately.
\end{proof}


\section{Regularization properties of \texorpdfstring{\Cref{alg:error-iter,alg:mixed-iter}}{the error-equation and mixed iterative schemes}}\label{sec:convergence_error}

We now analyze the error-equation and mixed iterations. Throughout this section, the graph regularizer is still
\(\Rreg_\bz^q\) from \Cref{def:graph-regularizer}, and the solution operator
\(T_{\alpha,\bz}^q\) is the one defined in
\Cref{eq:fixed-graph-functional}. As in \Cref{sec:iterschemes}, when
\(T_{\alpha,\bz}^q\) is set-valued, one admissible minimizer is selected.

For \Cref{alg:error-iter}, we use the residual convention
\[
\br_n^\delta \coloneqq \by^\delta-K\bx_n^\delta .
\]
Then the error-equation scheme can be written compactly as
\begin{equation}\label{eq:error-scheme-shifted}
\left\{
\begin{aligned}
\br_n^\delta &= \by^\delta-K\bx_n^\delta,\\
\bh_{n+1}^\delta
&\in T_{\alpha_{n+1},\bh_n^\delta}^q(\br_n^\delta),\\
\bx_{n+1}^\delta
&=\bx_n^\delta+\bh_{n+1}^\delta,
\end{aligned}
\right.
\qquad n\ge0,
\end{equation}
where \(\bh_0^\delta\in X\) is the prescribed image-space initialization used to
construct the first error graph. Equivalently,
\[
\br_{n+1}^\delta
=
\br_n^\delta-K\bh_{n+1}^\delta
=
\by^\delta-K\bx_{n+1}^\delta .
\]

\subsection{One correction step}\label{sec:error-one-step}

For a residual \(\br\in Y\), a reference \(\bz\in X\), and
\(\alpha>0\), the correction step is
\[
\bh\in T_{\alpha,\bz}^q(\br),
\]
that is,
\[
\bh\in
\operatorname*{argmin}_{\bu\in X}
\left\{
\frac12\norm{K\bu-\br}_Y^2+\alpha\Rreg_\bz^q(\bu)
\right\}.
\]
This is exactly the fixed-reference problem
\(\Phi_{\alpha,\bz}^q(\cdot;\br)\) from
\Cref{eq:fixed-graph-functional}.

\begin{proposition}\label{prop:error-one-step}
Assume \Cref{hyp:joint-coercivity}. Let \(q\in[1,2]\). For every
\(\br\in Y\), every \(\bz\in X\), and every \(\alpha>0\), the set
\(T_{\alpha,\bz}^q(\br)\) is nonempty. More precisely, if
\[
\Phi_{\alpha,\bz}^q(\bh;\br)\le M,
\]
then
\begin{equation}\label{eq:error-one-step-sublevel}
\norm{\bh}_2
\le
C_J\left(
(2M)^{1/2}+\norm{\br}_Y
+\left(\frac{qM}{\alpha}\right)^{1/q}
\right).
\end{equation}
The bound is independent of the reference image \(\bz\) and of the
discretization level \(N\). If \(q>1\), then \(T_{\alpha,\bz}^q(\br)\) is a
singleton.
\end{proposition}

\begin{proof}
This is \Cref{prop:existence} applied with data \(\by=\br\) and minimization
variable \(\bx=\bh\). The estimate \cref{eq:error-one-step-sublevel} is exactly
\cref{eq:sublevel-bound} with this notation.
\end{proof}

\begin{proposition}\label{prop:error-one-step-stability}
Assume \Cref{hyp:joint-coercivity}. Fix \(q\in[1,2]\) and \(\alpha>0\). Let
\[
\br_m\to\br \quad\text{in }Y,
\qquad
\norm{\bz_m-\bz}_\infty\to0,
\]
and choose
\[
\bh_m\in T_{\alpha,\bz_m}^q(\br_m).
\]
Then \((\bh_m)\) is bounded. Every cluster point of \((\bh_m)\) belongs to
\(T_{\alpha,\bz}^q(\br)\). If the limiting minimizer is unique, in particular if
\(q>1\), then the whole sequence converges to that minimizer.
\end{proposition}

\begin{proof}
This is \Cref{cor:one-step-stability} with \(\by_m=\br_m\),
\(\by=\br\), and \(\bx_m=\bh_m\).
\end{proof}

\subsection{Finite-iteration stability of \texorpdfstring{\Cref{alg:error-iter}}{the error-equation scheme}}\label{sec:error-finite-stability}
As for the standard scheme, one-step stability propagates by induction through
any fixed number of iterations of the error-equation scheme
\cref{eq:error-scheme-shifted}. The difference is that now we need to track the pair
$(\bx_n,\bh_n)$ of current reconstruction and graph input reference, so convergence is
required not only of the data and the initial reconstruction but also of the
initial graph input $\bh_0$. Under these assumptions the first $J$ iterates
converge as well, with bounds independent of the discretization level $N$.
\begin{theorem}\label{thm:error-finite-stability}
Assume \Cref{hyp:joint-coercivity}. Fix \(J\in\N\) and positive parameters
\(\alpha_1,\ldots,\alpha_J\). Let
\[
\by_m\to \by \quad\text{in }Y,
\qquad
\bx_{0,m}\to \bx_0 \quad\text{in }\norm{\cdot}_2,
\qquad
\norm{\bh_{0,m}-\bh_0}_\infty\to0 .
\]
For each \(m\), let \((\bx_{n,m},\bh_{n,m})\) be generated by
\Cref{eq:error-scheme-shifted} with data \(\by_m\), initial reconstruction
\(\bx_{0,m}\), and initial graph input \(\bh_{0,m}\), for
\(n=0,\ldots,J\).

Then, at each fixed discretization level, the family
\[
\{\bx_{n,m}:0\le n\le J\}\cup\{\bh_{n,m}:0\le n\le J\}
\]
is bounded. If the data and the initial families are bounded by constants
independent of \(N\), then the same is true for the generated finite family.

Moreover, every subsequence has a further subsequence such that, for each
\(n=0,\ldots,J\),
\[
\bx_{n,m_k}\to \bar\bx_n \quad\text{in } X,
\qquad
\norm{\bh_{n,m_k}-\bar\bh_n}_\infty\to0,
\]
where \((\bar\bx_n,\bar\bh_n)_{n=0}^J\) is generated by
\Cref{eq:error-scheme-shifted} with limiting data \(\by\), initial reconstruction
\(\bx_0\), and initial graph input \(\bh_0\). Equivalently,
\[
\bar\br_n=\by-K\bar\bx_n,
\qquad
\bar\bh_{n+1}\in T_{\alpha_{n+1},\bar\bh_n}^q(\bar\br_n),
\qquad
\bar\bx_{n+1}=\bar\bx_n+\bar\bh_{n+1}.
\]
If the minimizer is unique at every limiting correction step, in particular if
\(q>1\), then the whole sequence converges at each fixed step.
\end{theorem}

\begin{proof}
We argue by induction. The convergence of \(\bx_{0,m}\) and \(\bh_{0,m}\) is
assumed. Suppose that, along a subsequence,
\[
\bx_{n,m}\to \bar\bx_n,
\qquad
\norm{\bh_{n,m}-\bar\bh_n}_\infty\to0 .
\]
Then
\[
\br_{n,m}\coloneqq \by_m-K\bx_{n,m}
\to
\bar\br_n\coloneqq \by-K\bar\bx_n
\quad\text{in }Y .
\]
Since
\[
\bh_{n+1,m}\in T_{\alpha_{n+1},\bh_{n,m}}^q(\br_{n,m}),
\]
\Cref{prop:error-one-step-stability} gives boundedness of
\((\bh_{n+1,m})\), and every cluster point solves the limiting correction
problem
\[
\bar\bh_{n+1}\in T_{\alpha_{n+1},\bar\bh_n}^q(\bar\br_n).
\]
After extracting a further subsequence,
\[
\bh_{n+1,m}\to \bar\bh_{n+1}
\]
in \(X\), hence also in \(\norm{\cdot}_\infty\) at fixed \(N\). The update
\[
\bx_{n+1,m}=\bx_{n,m}+\bh_{n+1,m}
\]
then gives
\[
\bx_{n+1,m}\to \bar\bx_n+\bar\bh_{n+1}\eqqcolon \bar\bx_{n+1}.
\]
This proves the induction step. If each limiting correction step has a unique
minimizer, no subsequence extraction is needed, and the whole sequence converges.
\end{proof}

\subsection{Convergence of \texorpdfstring{\Cref{alg:error-iter}}{the error-equation scheme} for noisy data}
\label{sec:error-noisy}

In contrast to \Cref{alg:standard-iter}, the regularization term in
\Cref{alg:error-iter} acts on the correction \(\bh_{n+1}^\delta\), not directly
on the cumulative reconstruction \(\bx_n^\delta\). Consequently, without an
additional selection principle, the natural convergence statement is convergence
of the stopped reconstructions to the exact-solution set \(\Ssol(\by)\), rather
than to the graph-minimizing set \(\mathcal S_{\rm gm}(\by)\).

Let \(m(\delta)\ge1\) be an a priori stopping index satisfying
\begin{equation}\label{eq:error-apriori-rule}
\alpha_{m(\delta)}\to0,
\qquad
\frac{\delta^2}{\alpha_{m(\delta)}}\to0
\qquad\text{as }\delta\downarrow0 .
\end{equation}
We assume that the predecessor of the stopped reconstruction is uniformly
bounded:
\begin{equation}\label{eq:error-noisy-bounded-x}
\sup_{0<\delta\le\delta_0}
\norm{\bx_{m(\delta)-1}^\delta}_2
\le M_X,
\end{equation}
with \(M_X\) independent of \(N\). 

\begin{theorem}
\label{thm:error-noisy-convergence}
Assume \Cref{hyp:consistency,hyp:joint-coercivity}. Let
\(\norm{\by^\delta-\by}_Y\le\delta\), and let
\((\bx_n^\delta,\bh_n^\delta)\) be generated by
\Cref{eq:error-scheme-shifted}. Suppose that
\Cref{eq:error-apriori-rule,eq:error-noisy-bounded-x} hold. Then the stopped
family \((\bx_{m(\delta)}^\delta)\) is bounded in \(\norm{\cdot}_2\), and
\begin{equation}\label{eq:error-residual-convergence}
\norm{K\bx_{m(\delta)}^\delta-\by}_Y\to0
\qquad\text{as }\delta\downarrow0.
\end{equation}
Moreover,
\begin{equation}\label{eq:error-set-convergence}
\operatorname{dist}_2
\bigl(\bx_{m(\delta)}^\delta,\Ssol(\by)\bigr)
\to0
\qquad\text{as }\delta\downarrow0.
\end{equation}
\end{theorem}

\begin{proof}
Set \(m=m(\delta)\). Let \(\gt\) be the bounded exact solution from
\Cref{hyp:consistency}, and set
\[
\be_{m-1}^\delta\coloneqq \gt-\bx_{m-1}^\delta .
\]
Since \(K\gt=\by\) and
\(\br_{m-1}^\delta=\by^\delta-K\bx_{m-1}^\delta\),
\[
K\be_{m-1}^\delta-\br_{m-1}^\delta
=
\by-\by^\delta,
\qquad
\norm{K\be_{m-1}^\delta-\br_{m-1}^\delta}_Y\le\delta .
\]
Moreover, because \(q\le2\),
\[
\norm{\be_{m-1}^\delta}_q
\le
\norm{\gt}_q+\norm{\bx_{m-1}^\delta}_q
\le
M_{\rm gt}+M_X .
\]
Hence, by \Cref{eq:R-uniform-bound},
\[
\Rreg_{\bh_{m-1}^\delta}^q(\be_{m-1}^\delta)
\le
C_E,
\qquad
C_E\coloneqq
\frac{(2d_R)^q}{q}(M_{\rm gt}+M_X)^q .
\]
Minimality of \(\bh_m^\delta\) gives
\[
\frac12\norm{K\bh_m^\delta-\br_{m-1}^\delta}_Y^2
+
\alpha_m\Rreg_{\bh_{m-1}^\delta}^q(\bh_m^\delta)
\le
\frac12\delta^2+\alpha_m C_E .
\]
Since
\[
K\bh_m^\delta-\br_{m-1}^\delta
=
-\br_m^\delta,
\]
we obtain
\begin{equation*}
\norm{\br_m^\delta}_Y^2
\le
\delta^2+2\alpha_m C_E .
\end{equation*}
By \Cref{eq:error-apriori-rule}, the right-hand side tends to zero. Therefore
\[
\norm{K\bx_m^\delta-\by}_Y
\le
\norm{K\bx_m^\delta-\by^\delta}_Y
+\norm{\by^\delta-\by}_Y
=
\norm{\br_m^\delta}_Y
+\norm{\by^\delta-\by}_Y
\to0,
\]
which proves \Cref{eq:error-residual-convergence}.

It remains to show the boundedness of the stopped reconstructions. Define
\[
M_\delta\coloneqq \frac12\delta^2+\alpha_m C_E .
\]
The bound in \Cref{eq:error-noisy-bounded-x} and \Cref{hyp:joint-coercivity} give
\[
\norm{\br_{m-1}^\delta}_Y
\le
\norm{\by}_Y+\delta_0+M_KM_X
\]
for \(0<\delta\le\delta_0\). Applying the sublevel estimate
\Cref{eq:error-one-step-sublevel} to \(\bh_m^\delta\) yields
\[
\norm{\bh_m^\delta}_2
\le
C_J\left(
(2M_\delta)^{1/2}
+\norm{\br_{m-1}^\delta}_Y
+\left(\frac{qM_\delta}{\alpha_m}\right)^{1/q}
\right).
\]
The first term is bounded, the second term is bounded by the previous estimate,
and
\[
\frac{M_\delta}{\alpha_m}
=
\frac{\delta^2}{2\alpha_m}+C_E
\]
is bounded for small \(\delta\) by \Cref{eq:error-apriori-rule}. Hence
\((\bh_m^\delta)\) is bounded in \(\norm{\cdot}_2\). Together with
\Cref{eq:error-noisy-bounded-x}, this gives boundedness of
\[
\bx_m^\delta=\bx_{m-1}^\delta+\bh_m^\delta .
\]

We now prove \Cref{eq:error-set-convergence}. Suppose, by contradiction, that it
fails. Then there exist \(\varepsilon>0\) and a sequence
\(\delta_j\downarrow0\) such that
\[
\operatorname{dist}_2
\bigl(\bx_{m(\delta_j)}^{\delta_j},\Ssol(\by)\bigr)
\ge\varepsilon
\qquad\forall j.
\]
By the boundedness just proved and finite-dimensional compactness, after passing
to a subsequence,
\[
\bx_{m(\delta_j)}^{\delta_j}\to\bx^\star .
\]
The residual convergence \Cref{eq:error-residual-convergence} gives
\(K\bx^\star=\by\), hence \(\bx^\star\in\Ssol(\by)\). Therefore
\[
\operatorname{dist}_2
\bigl(\bx_{m(\delta_j)}^{\delta_j},\Ssol(\by)\bigr)
\le
\norm{\bx_{m(\delta_j)}^{\delta_j}-\bx^\star}_2
\to0,
\]
contradicting the choice of the sequence. This proves
\Cref{eq:error-set-convergence}.
\end{proof}

\begin{remark}
The boundedness condition \cref{eq:error-noisy-bounded-x} is the only additional condition required by the error-equation scheme but not by the standard scheme. In
\Cref{alg:standard-iter}, joint coercivity controls the same variable that is
being reconstructed, because the functional contains
\(\Rreg_{\bx_{n-1}^\delta}^q(\bx_n^\delta)\). In
\Cref{alg:error-iter}, the functional controls the correction
\(\bh_m^\delta\), while the reconstruction is the cumulative quantity
\(\bx_m^\delta=\bx_{m-1}^\delta+\bh_m^\delta\). Thus joint coercivity alone does
not control the previous reconstruction \(\bx_{m-1}^\delta\); some boundedness or stopping
selection condition is needed unless additional structure is imposed. However, we remark that for grayscale images, such a bound is easily verified.
\end{remark}

\subsection{Convergence of \texorpdfstring{\Cref{alg:mixed-iter}}{the mixed scheme} for noisy data}
\label{sec:mixed-convergence}

The mixed scheme consists of a fixed standard phase followed by an error-equation
phase. Since the standard phase has fixed length, its finite-iteration stability
follows from \Cref{thm:finite-stability}; after that point, the error phase is
exactly \Cref{alg:error-iter} with shifted indices. Therefore the vanishing-noise
argument is inherited from \Cref{thm:error-noisy-convergence}.

Fix \(n_{\mathrm{std}}\in\mathbb N_{>0}\). After the standard phase, write the
chosen initial graph input for the error stage as \(\bg_0^\delta\), and set
\[
\br_{n_{\mathrm{std}}}^\delta
=
\by^\delta-K\bx_{n_{\mathrm{std}}}^\delta .
\]
For \(k\ge0\), the error stage of \Cref{alg:mixed-iter} is
\begin{equation}\label{eq:mixed-error-stage}
\left\{
\begin{aligned}
\br_{n_{\mathrm{std}}+k}^\delta
&=
\by^\delta-K\bx_{n_{\mathrm{std}}+k}^\delta,\\
\bg_{k+1}^\delta
&\in
T_{\alpha_{n_{\mathrm{std}}+k+1},\bg_k^\delta}^q
\bigl(\br_{n_{\mathrm{std}}+k}^\delta\bigr),\\
\bx_{n_{\mathrm{std}}+k+1}^\delta
&=
\bx_{n_{\mathrm{std}}+k}^\delta+\bg_{k+1}^\delta .
\end{aligned}
\right.
\end{equation}

Let the number \(m(\delta)\ge1\) of error-stage iterations satisfy
\begin{equation}\label{eq:mixed-apriori-rule}
\alpha_{n_{\mathrm{std}}+m(\delta)}\to0,
\qquad
\frac{\delta^2}{\alpha_{n_{\mathrm{std}}+m(\delta)}}\to0
\qquad\text{as }\delta\downarrow0 .
\end{equation}
Assume that the predecessor of the stopped mixed reconstruction is uniformly
bounded:
\begin{equation}\label{eq:mixed-noisy-bounded}
\sup_{0<\delta\le\delta_0}
\norm{\bx_{n_{\mathrm{std}}+m(\delta)-1}^\delta}_2
\le M_X,
\end{equation}
with \(M_X\) independent of \(N\).

\begin{theorem}
\label{thm:mixed-noisy-convergence}
Assume \Cref{hyp:consistency,hyp:joint-coercivity}. Let
\(\norm{\by^\delta-\by}_Y\le\delta\). For each \(\delta\), compute
\(n_{\mathrm{std}}\) standard iterations and then apply the error stage
\Cref{eq:mixed-error-stage}. Suppose that
\Cref{eq:mixed-apriori-rule,eq:mixed-noisy-bounded} hold. Then the stopped mixed
family
\[
\bigl(\bx_{n_{\mathrm{std}}+m(\delta)}^\delta\bigr)
\]
is bounded in \(\norm{\cdot}_2\), and
\begin{equation}\label{eq:mixed-residual-convergence}
\norm{K\bx_{n_{\mathrm{std}}+m(\delta)}^\delta-\by}_Y\to0
\qquad\text{as }\delta\downarrow0.
\end{equation}
Moreover,
\begin{equation}\label{eq:mixed-set-convergence}
\operatorname{dist}_2
\bigl(\bx_{n_{\mathrm{std}}+m(\delta)}^\delta,\Ssol(\by)\bigr)
\to0
\qquad\text{as }\delta\downarrow0 .
\end{equation}
\end{theorem}

\begin{proof}
Apply \Cref{thm:error-noisy-convergence} to the shifted variables
\[
\widetilde\bx_k^\delta
\coloneqq
\bx_{n_{\mathrm{std}}+k}^\delta,
\qquad
\widetilde\bh_k^\delta
\coloneqq
\bg_k^\delta,
\qquad
\widetilde\alpha_k
\coloneqq
\alpha_{n_{\mathrm{std}}+k}.
\]
Then \Cref{eq:mixed-error-stage} is exactly
\Cref{eq:error-scheme-shifted} for
\((\widetilde\bx_k^\delta,\widetilde\bh_k^\delta)\). The parameter rule
\Cref{eq:mixed-apriori-rule} is \Cref{eq:error-apriori-rule} for
\(\widetilde\alpha_{m(\delta)}\), and the boundedness condition
\Cref{eq:mixed-noisy-bounded} is \Cref{eq:error-noisy-bounded-x} for the shifted
predecessor
\[
\widetilde\bx_{m(\delta)-1}^\delta
=
\bx_{n_{\mathrm{std}}+m(\delta)-1}^\delta .
\]
Therefore \Cref{thm:error-noisy-convergence} gives boundedness of
\[
\widetilde\bx_{m(\delta)}^\delta
=
\bx_{n_{\mathrm{std}}+m(\delta)}^\delta,
\]
as well as \Cref{eq:mixed-residual-convergence} and
\Cref{eq:mixed-set-convergence}.
\end{proof}

\begin{remark}
The fixed standard phase does not need to converge to an exact solution for
\Cref{thm:mixed-noisy-convergence}; it only supplies the initial reconstruction
and initial residual for the shifted error-equation scheme. 
\end{remark}

\section{Numerical tests}\label{sec:numerical_results}

In this section, we present numerical results for image deblurring and CT problems
obtained with the proposed iterative schemes described in \Cref{sec:iterschemes}.
We consider $\ell^2$--$\ell^1$ regularization, i.e., $q=1$. The approximate solution of the resulting $\ell^2$--$\ell^1$ minimization problems is computed by the majorization-minimization generalized Krylov subspace (MM-GKS) method with the adaptive majorant proposed in \cite{LMRS15} as implemented in \cite{buccini2024software}.
This implementation employs a restarting strategy for the generalized Krylov subspace to reduce the computational cost, while the regularization parameter is automatically estimated by GCV.



For every fixed deblurring or CT matrix used below,
\Cref{rem:finite-dim-coercivity} gives the joint coercivity needed for the
$q=1$ numerical problem, although its constant may depend on $N$. At the
dimension-free level, \Cref{prop:dim-free-deblurring} treats the periodic blur
model for $q=2$, whereas \Cref{prop:dim-free-CT} treats the full-angle CT model
for every $q\in[1,2]$.


The parameters defining the graph Laplacian are set to $R=5$ and $\sigma=10^{-3}$ for the standard iterative scheme \cref{eq:standard-iter}; for the error-equation stage of the mixed iterative scheme in \Cref{alg:mixed-iter}, we keep $R=5$ and set $\sigma=10^{-4}$.
For the finite computations reported here, the contrast cap $B_{\rm c}$ in \Cref{eq:contrast-cap} is chosen larger
than $1$, so the cap
is inactive and the computed Gaussian weights are unchanged. 


We compare the reconstruction quality of our iterated graph Laplacian methods with that obtained after the first iteration, i.e., the graphLa\texttt{+}$\Psi$ method proposed in \cite{bianchi2025data}.
To assess the quality of the reconstructions, we consider the standard imaging performance metrics: the Relative Reconstruction Error (RRE), the Peak Signal-to-Noise Ratio (PSNR), the Structural Similarity Index Measure (SSIM), and the Gradient Magnitude Similarity Deviation (GMSD) \cite{Xue2014-wt}, which is closely related to SSIM and assesses the similarity between two images through the standard deviation of the gradient magnitude similarity map.
When we present the performance metrics of a nonrectangular region of interest (ROI), SSIM and GMSD are computed on the smallest rectangular subimage that contains the ROI.

In the tables below, rows above the horizontal line report standard iterates $\bx_n^\delta$, while rows below the horizontal line report mixed reconstructions $\bx_{n_{\mathrm{std}}+k}^\delta$ after $k$ error-equation steps, consistent with \Cref{eq:mixed-error-stage}.


\subsection[Test 1: Satellite deblurring with Tikhonov initialization]{Test 1: Satellite deblurring for $\bx_0^\delta$ computed by $\ell^2$--$\ell^2$ Tikhonov}
For the first image deblurring problem, we consider the images in \Cref{fig:test4}.
The test image is a $238 \times 238$ grayscale image, and the PSF has a support of $9 \times 9$ pixels. Gaussian noise at level $1\%$ has been added to the blurred image to obtain the observed corrupted image.

\begin{figure}
        \centering
		\begin{subfigure}{0.25\textwidth}
			\includegraphics[width=\textwidth]{figures/Numerical_Tests/test_satellite/test5GCV_TRUE}
			\caption{}
		\end{subfigure}\qquad
		\begin{subfigure}{0.25\textwidth}
			\includegraphics[clip,trim={2cm 1.5cm 2cm 1.5cm}, width=\textwidth, height=\textwidth]{figures/Numerical_Tests/test_satellite/test5GCV_PSF}
			\caption{}
		\end{subfigure}\qquad
		\begin{subfigure}{0.25\textwidth}
			\includegraphics[width=\textwidth]{figures/Numerical_Tests/test_satellite/test5GCV_CORR}
			\caption{}
		\end{subfigure}
		\caption{Satellite deblurring test: (a) original image, (b) PSF with a $9\times9$-pixel support, (c) corrupted image.}
		\label{fig:test4}
	\end{figure}

The initial approximation $\bx_0^\delta$ is computed by $\ell^2$--$\ell^2$ Tikhonov regularization.
Note that the solution computed after the first iteration, i.e., $\bx_1^\delta$, is exactly the one computed by the graphLa\texttt{+}Tik method proposed in \cite{bianchi2025data}.

For the mixed iteration, in the notation of \Cref{eq:mixed-error-stage}, the initial error-stage graph input is $\bg_0^\delta=\br_{n_{\mathrm{std}}}^\delta$, consistent with the choice in \Cref{eq:g0res}.
The results for the general approach, which computes $\bg_0^\delta$ using Tikhonov regularization to approximate the solution of the error equation $K\bg=\br_{n_{\mathrm{std}}}^\delta$, are not reported since they are indistinguishable from the simple choice $\bg_0^\delta=\br_{n_{\mathrm{std}}}^\delta$. 

		\begin{table}
    \centering
        \begin{tabular}{c|c|c|c | c }
            iterate & RRE & PSNR & SSIM & GMSD \\ \hline
                $\bx_{0}^{\delta}$ & 0.2178  & 26.23 &  0.8949 & 0.0739\\
			$\bx_{1}^{\delta}$  & 0.1434 & 29.87 & 0.9565 & 0.0290 \\ 
			$\bx_{2}^{\delta}$  & 0.1222 & 31.25 & 0.9679 & 0.0209 \\ 
			$\bx_{3}^{\delta}$  & 0.1155 & 31.74 & 0.9721 & 0.0197 \\ 
			$\bx_{4}^{\delta}$  & 0.1109 & 32.10 & 0.9744 & 0.0193  \\ 
			$\bx_{n_{\mathrm{std}}}^{\delta}$  & 0.1099 & 32.17 & 0.9756 & 0.0194 \\ \hline
            $\bx_{n_{\mathrm{std}}+1}^{\delta}$  & 0.1093 & 32.22 & 0.9737 & 0.0176\\ 
			$\bx_{n_{\mathrm{std}}+2}^{\delta}$  & 0.1089 & 32.25 & 0.9713 & 0.0165\\ 
			$\bx_{n_{\mathrm{std}}+3}^{\delta}$  & 0.1086 & 32.28 & 0.9684 & 0.0157\\ 
			$\bx_{n_{\mathrm{std}}+4}^{\delta}$  & 0.1084 & 32.29 & 0.9651 & 0.0151\\ 
			$\bx_{n_{\mathrm{std}}+5}^{\delta}$  & 0.1083 & 32.30 & 0.9614 & 0.0146\\ 
			$\bx_{n_{\mathrm{std}}+6}^{\delta}$  & 0.1083 & 32.30 & 0.9575 & 0.0143\\ 
        \end{tabular}
		\caption{Performance metrics for Test 1. The upper part is the standard iterative scheme with $\bx_{0}^{\delta}$ computed by Tikhonov. The lower part is the error-equation stage of the mixed iterative scheme, started from $\bx_{n_{\mathrm{std}}}^{\delta}=\bx_{5}^{\delta}$.}
		\label{Table: test4}
	\end{table}
    
The performance metrics reported in \Cref{Table: test4} confirm the significant improvement achieved after the first iteration, in agreement with the results in~\cite{bianchi2025data}. 
In particular, the GMSD values confirm the effectiveness of the error-equation stage of the mixed iterative scheme.

The improvement in the quality of the reconstructions is more evident when looking at the restored images in \Cref{fig:test4 rec}, where the reconstruction for the graphLa\texttt{+}Tik ($\bx_{1}^{\delta}$) is compared with $\bx_{n_{\mathrm{std}}}^{\delta}$, and with the final reconstruction computed by the mixed iterative scheme. We note an improvement in the details, which is even more evident in \Cref{fig:test4 detail 1}, where a detail in the red box is considerably improved by the mixed iterative scheme, as confirmed by the metrics reported in the caption.

	\begin{figure}
    \centering
		\begin{subfigure}{0.25\textwidth}
			\includegraphics[width=\textwidth]{figures/Numerical_Tests/test_satellite/test5GCV_graphLa}
			\caption{}
		\end{subfigure}\qquad
		\begin{subfigure}{0.25\textwidth}
			\includegraphics[width=\textwidth]{figures/Numerical_Tests/test_satellite/test5GCV_rec_iter}
			\caption{}
		\end{subfigure}\qquad
		\begin{subfigure}{0.25\textwidth}
			\includegraphics[width=\textwidth]{figures/Numerical_Tests/test_satellite/test5GCV_rec_res_no_st_6}
			\caption{}
		\end{subfigure}
		\caption{Restored images for Test 1: (a) graphLa\texttt{+}Tik ($\bx_{1}^{\delta}$), (b) standard iterative scheme ($\bx_{n_{\mathrm{std}}}^{\delta}=\bx_{5}^{\delta}$), (c) final reconstruction with mixed iterative scheme ($\bx_{n_{\mathrm{std}}+6}^{\delta}$).}
		\label{fig:test4 rec}
	\end{figure}
	
	\begin{figure}
    	\begin{subfigure}{0.24\textwidth}
			\includegraphics[width=\textwidth]{figures/Numerical_Tests/test_satellite/test5GCV_ROI_true.eps}
			\caption{}
		\end{subfigure}
		\begin{subfigure}{0.24\textwidth}
			\includegraphics[width=\textwidth]{figures/Numerical_Tests/test_satellite/test5GCV_detail1_graphLa}
			\caption{}
		\end{subfigure}
		\begin{subfigure}{0.24\textwidth}
			\includegraphics[width=\textwidth]{figures/Numerical_Tests/test_satellite/test5GCV_detail1_rec_iter}
			\caption{}
		\end{subfigure}
		\begin{subfigure}{0.24\textwidth}
			\includegraphics[width=\textwidth]{figures/Numerical_Tests/test_satellite/test5GCV_detail1_rec_res_no_st_6}
			\caption{}
		\end{subfigure}
        \caption{Details of the restored images for Test 1: (a) detail of $\gt$ in the red box for the computed metrics, (b) graphLa\texttt{+}Tik (PSNR = 19.3937, SSIM = 0.6933, GMSD = 0.0933), (c) standard iterative scheme (PSNR = 20.6415, SSIM = 0.7836, GMSD = 0.0820), (d) final reconstruction (PSNR = 21.1788, SSIM = 0.8074, GMSD = 0.0583).}
		\label{fig:test4 detail 1}
	\end{figure}


\subsection[Test 2: Satellite deblurring with TV initialization]{Test 2: Satellite deblurring for $\bx_0^\delta$ computed by TV}

In this example, we show the importance of a good initial approximation $\bx_0^\delta$.
It is obtained via a total variation (TV) reconstruction, corresponding to the minimization problem \cref{model_eq2} with $q=1$ and a finite-difference approximation of the gradient as the operator $D$.
The solution computed after the first iteration is referred to as the graphLa\texttt{+}TV reconstruction, following the notation in \cite{bianchi2025data}.

    We consider again the previous example in \Cref{fig:test4}, using the new initial approximation \(\bx_0^\delta\) computed via TV regularization. The performance metrics reported in \Cref{Table: test_TV} show an improvement compared with those in \Cref{Table: test4}. In this case, the additional improvement obtained through further iterations is less pronounced. Nevertheless, an enhancement of the image details can still be observed, as evidenced by the reconstructed images in \Cref{fig:test_TV rec,fig:test_TV detail 1}. The performance metrics in \Cref{fig:test_TV detail 1} refer to the red-boxed subimage, as in \Cref{fig:test4 detail 1}.

    \begin{table}
    \centering
        \begin{tabular}{c|c|c|c | c}
           iterate & RRE & PSNR & SSIM & GMSD \\ \hline
                $\bx_{0}^{\delta}$ & 0.1418  & 29.96 & 0.9077 & 0.0278\\
                $\bx_{1}^{\delta}$  & 0.1076 & 32.36 & 0.9759 & 0.0124\\ 
			$\bx_{2}^{\delta}$  & 0.0975 & 33.22 & 0.9814 & 0.0096 \\ 
			$\bx_{3}^{\delta}$  & 0.0941 & 33.53 & 0.9832 & 0.0093 \\ 
		$\bx_{n_{\mathrm{std}}}^{\delta}$  & 0.0939 & 33.54 & 0.9834 & 0.0091 \\ \hline
        $\bx_{n_{\mathrm{std}}+1}^{\delta}$  &  0.0937 & 33.56 & 0.9811 & 0.0087\\ 
		$\bx_{n_{\mathrm{std}}+2}^{\delta}$  &  0.0936 & 33.57 & 0.9783 & 0.0085\\ 
        \end{tabular}
		\caption{Performance metrics for Test 2. The upper part is the standard iterative scheme with $\bx_{0}^{\delta}$ computed by TV. The lower part is the error-equation stage of the mixed iterative scheme, started from $\bx_{n_{\mathrm{std}}}^{\delta}=\bx_{4}^{\delta}$.}
		\label{Table: test_TV}
	\end{table}

\begin{figure}
        \centering
		\begin{subfigure}{0.25\textwidth}
			\includegraphics[width=\textwidth]{figures/Numerical_Tests/test_satellite_TV/test_TV_graphLa}
			\caption{}
		\end{subfigure}\qquad
		\begin{subfigure}{0.25\textwidth}
			\includegraphics[width=\textwidth]{figures/Numerical_Tests/test_satellite_TV/test_TV_rec_iter}
			\caption{}
		\end{subfigure}\qquad
		\begin{subfigure}{0.25\textwidth}
			\includegraphics[width=\textwidth]{figures/Numerical_Tests/test_satellite_TV/test_TV_rec_res}
			\caption{}
		\end{subfigure}
		\caption{Restored images for Test 2: (a) graphLa\texttt{+}TV ($\bx_{1}^{\delta}$), (b) standard iterative scheme ($\bx_{n_{\mathrm{std}}}^{\delta}=\bx_{4}^{\delta}$), (c) final reconstruction with mixed iterative scheme ($\bx_{n_{\mathrm{std}}+2}^{\delta}$).}
		\label{fig:test_TV rec}
	\end{figure}
	
	\begin{figure}
        \centering
        \begin{subfigure}{0.24\textwidth}
			\includegraphics[width=\textwidth]{figures/Numerical_Tests/test_satellite_TV/test_TV_ROI_true.eps}
			\caption{}
		\end{subfigure}
		\begin{subfigure}{0.24\textwidth}
			\includegraphics[width=\textwidth]{figures/Numerical_Tests/test_satellite_TV/test_TV_detail1_graphLa}
			\caption{}
		\end{subfigure}
		\begin{subfigure}{0.24\textwidth}
			\includegraphics[width=\textwidth]{figures/Numerical_Tests/test_satellite_TV/test_TV_detail1_rec_iter}
			\caption{}
		\end{subfigure}
		\begin{subfigure}{0.24\textwidth}
			\includegraphics[width=\textwidth]{figures/Numerical_Tests/test_satellite_TV/test_TV_detail1_rec_res}
			\caption{}
		\end{subfigure}
        \caption{Details of the restored images for Test 2: (a) detail of $\gt$ in the red box for the computed metrics, (b) graphLa\texttt{+}TV (PSNR = 20.9493, SSIM = 0.8014, GMSD = 0.0610), (c) standard iterative scheme (PSNR = 21.5745, SSIM = 0.8459, GMSD = 0.0448), (d) final reconstruction (PSNR = 21.6567, SSIM = 0.8473, GMSD = 0.0435).}
		\label{fig:test_TV detail 1}
	\end{figure}


\subsection{Test CT: a computed tomography example}\label{ssec:CT}

In this numerical test, the image is a $256 \times 256$ grayscale image from a real CT scan, taken from the AAPM Low Dose CT Grand Challenge dataset~\cite{Moen2021-me}. The sinogram is computed using AIRToolsII~\cite{Hansen2017-bb}. We consider a simulated CT acquisition with 60 equispaced projection angles in the interval $[0,\pi]$ and 362 detector bins. To simulate real-world scenarios, the sinogram is corrupted with additive Gaussian noise at a noise level of $1\%$. The original image and the noisy sinogram are shown in \Cref{fig:test_CT_REAL}.

The parameter $\sigma$ used for the construction of the graph Laplacian is set to $\sigma=5\times10^{-3}$.
We apply the standard iterative scheme \cref{eq:standard-iter}, with the initial approximation $\bx_{0}^{\delta}$ computed via FBP. The standard iterative scheme is performed for two iterations. Subsequently, one error-equation step in the mixed scheme is applied, where the Laplacian is constructed using the back-projected residual obtained by FBP with Hann filtering.

The performance metrics reported in \Cref{Table: test_CT_REAL}, together with the restored images shown in \Cref{fig:test_CT_REAL rec}, support the effectiveness of the proposed iterative schemes. In particular, the close-up views in \Cref{fig:test_CT_REAL detail 1} highlight the ability of the error-equation stage to attenuate noise while preserving relevant structural details.
 
\begin{figure}
        \centering
		\begin{subfigure}{0.25\textwidth}
			\includegraphics[width=\textwidth]{figures/Numerical_Tests/CT_test/test_CT_true.eps}
			\caption{}
		\end{subfigure}\qquad
		\begin{subfigure}{0.25\textwidth}
			\includegraphics[clip,trim={2cm 1.5cm 2cm 1.5cm}, width=\textwidth, height=\textwidth]{figures/Numerical_Tests/CT_test/sinogram_REALeps.eps}
			\caption{}
		\end{subfigure}
		\caption{Test CT: (a) original image, (b) observed sinogram corrupted by noise.}
		\label{fig:test_CT_REAL}
	\end{figure}



\begin{table}
\centering
	\begin{tabular}{c|c|c|c | c}
       iterate & RRE & PSNR & SSIM & GMSD\\ \hline
		 $\bx_{0}^{\delta}$ & 0.3373 & 22.41 & 0.3389 & 0.1633 \\
    	$\bx_{1}^{\delta}$  & 0.1414 & 29.96 & 0.7532 & 0.0523 \\ 
		$\bx_{n_{\mathrm{std}}}^{\delta}$  & 0.1356 & 30.32 & 0.7828 & 0.0587 \\ \hline
		$\bx_{n_{\mathrm{std}}+1}^{\delta}$  & 0.1352 & 30.35 & 0.7644 & 0.0361 \\
	\end{tabular}
	\caption{Performance metrics for Test CT. The upper part is the standard iterative scheme with $\bx_{0}^{\delta}$ computed by FBP. Below the horizontal line is the error-equation stage of the mixed iterative scheme, started from $\bx_{n_{\mathrm{std}}}^{\delta}=\bx_{2}^{\delta}$.}
    \label{Table: test_CT_REAL}
\end{table}
    \begin{figure}
        \centering
		\begin{subfigure}{0.25\textwidth}
			\includegraphics[width=\textwidth]{figures/Numerical_Tests/CT_test_draft/GraphLa_sol}
			\caption{}
			\end{subfigure}\qquad
		\begin{subfigure}{0.25\textwidth}
			\includegraphics[width=\textwidth]{figures/Numerical_Tests/CT_test_draft/iter_sol}
			\caption{}			
			\end{subfigure}\qquad
		\begin{subfigure}{0.25\textwidth}
			\includegraphics[width=\textwidth]{figures/Numerical_Tests/CT_test_draft/res_sol}
			\caption{}
			\end{subfigure}
		\caption{Restored images for Test CT: (a) graphLa\texttt{+}FBP ($\bx_{1}^{\delta}$), (b) standard iterative scheme ($\bx_{n_{\mathrm{std}}}^{\delta}=\bx_{2}^{\delta}$), (c) final reconstruction with mixed iterative scheme ($\bx_{n_{\mathrm{std}}+1}^{\delta}$).}
        \label{fig:test_CT_REAL rec}
    \end{figure}

	    \begin{figure}
	        \centering
			\begin{subfigure}{0.24\textwidth}
	        \includegraphics[width=\textwidth]{figures/Numerical_Tests/CT_test_draft/new_roi}
				\caption{}
		\end{subfigure}
		\begin{subfigure}{0.24\textwidth}
			\includegraphics[width=\textwidth]{figures/Numerical_Tests/CT_test_draft/GraphLa_sol_detail}
			\caption{}
		\end{subfigure}
		\begin{subfigure}{0.24\textwidth}
			\includegraphics[width=\textwidth]{figures/Numerical_Tests/CT_test_draft/iter_sol_detail}
			\caption{}			
		\end{subfigure}
		\begin{subfigure}{0.24\textwidth}
			\includegraphics[width=\textwidth]{figures/Numerical_Tests/CT_test_draft/res_sol_detail}
			\caption{}
		\end{subfigure}
        \caption{Details of the restored images for Test CT: (a) detail of $\gt$ in the red box for the computed metrics, (b) graphLa\texttt{+}FBP (PSNR = 25.7133, SSIM = 0.7064, GMSD = 0.0590), (c) standard iterative scheme (PSNR = 25.7194, SSIM = 0.7361, GMSD = 0.0765), (d) final reconstruction (PSNR = 26.2945, SSIM = 0.7689, GMSD = 0.0337).}
        \label{fig:test_CT_REAL detail 1}
    \end{figure}


\subsection[Deblurring and denoising with neural network initialization]{Deblurring and denoising for $\bx_0^\delta$ computed by a neural network}

In this test, we use a deep Wiener convolution network \cite{dong2020deep} to compute the initial reconstruction $\bx_0^\delta$. We show that, even when using a state-of-the-art reconstruction method, the iterative graph Laplacian approach is still able to further improve the reconstruction quality.

The test images are grayscale images with intensity values scaled to the interval $[0,1]$. They are blurred using a $17\times17$ Gaussian kernel with standard deviation $5$, and additive Gaussian noise with standard deviation $0.01$ is subsequently added. For all iterations, we set $R=4$ and $\sigma=0.002$.

The mixed iterative scheme is then applied, consisting of $n_{\mathrm{std}}=2$ standard iterations followed by $3$ error-equation iterations.

The original image, the blurred image, $\bx_0^\delta$ (the initial approximation), $\bx_{n_{\mathrm{std}}}^{\delta}$ (the reconstruction after the standard iterative scheme), and $\bx_{n_{\mathrm{std}}+3}^{\delta}$ (the final reconstruction after the mixed scheme) are shown in \Cref{fig:network_test}. We also show details in a zoomed-in region in \Cref{fig:network_test_zoom}. The corresponding performance metrics are reported in \Cref{Table:network_test}.

The iterated graph Laplacian improves all four metrics (RRE, PSNR, SSIM, and GMSD) for both images. Although the mixed scheme provides only modest gains for the first image (House Front), it leads to a substantial improvement for the second image (Windows with Flowers). The benefits are particularly pronounced in the recovery of structural details, as shown in \Cref{fig:network_test_zoom} and reflected by the metrics reported in the caption.

\begin{figure}
    \centering
    
    \begin{subfigure}{0.19\textwidth}
    \includegraphics[width=\textwidth]{figures/Numerical_Tests/Network_Deblurring/org1.png}
    \includegraphics[width=\textwidth]{figures/Numerical_Tests/Network_Deblurring/org2.png}
    \caption{}
    \end{subfigure}
    \begin{subfigure}{0.19\textwidth}
    \includegraphics[width=\textwidth]{figures/Numerical_Tests/Network_Deblurring/blurred1.png}
    \includegraphics[width=\textwidth]{figures/Numerical_Tests/Network_Deblurring/blurred2.png}
    \caption{}
    \end{subfigure}
    \begin{subfigure}{0.19\textwidth}
    \includegraphics[width=\textwidth]{figures/Numerical_Tests/Network_Deblurring/deblurred1.png}
    \includegraphics[width=\textwidth]{figures/Numerical_Tests/Network_Deblurring/deblurred2.png}
    \caption{}
    \end{subfigure}
    \begin{subfigure}{0.19\textwidth}
    \includegraphics[width=\textwidth]{figures/Numerical_Tests/Network_Deblurring/graphLa1.png}
    \includegraphics[width=\textwidth]{figures/Numerical_Tests/Network_Deblurring/graphLa2.png}
    \caption{}
    \end{subfigure}
    \begin{subfigure}{0.19\textwidth}
    \includegraphics[width=\textwidth]{figures/Numerical_Tests/Network_Deblurring/graphLa_residual1.png}
    \includegraphics[width=\textwidth]{figures/Numerical_Tests/Network_Deblurring/graphLa_residual2.png}
    \caption{}
    \end{subfigure}
    \caption{Image deblurring and denoising with neural network initialization: (a) original image, (b) blurred and noisy image, (c) network reconstruction, (d) standard iterative scheme, (e) final reconstruction.}
    \label{fig:network_test}
\end{figure}

\begin{figure}
    \centering
    
    \begin{subfigure}{0.19\textwidth}
    \includegraphics[width=\textwidth]{figures/Numerical_Tests/Network_Deblurring/org1_zoom.png}
    \includegraphics[width=\textwidth]{figures/Numerical_Tests/Network_Deblurring/org2_zoom.png}
    \caption{}
    \end{subfigure}
    \begin{subfigure}{0.19\textwidth}
    \includegraphics[width=\textwidth]{figures/Numerical_Tests/Network_Deblurring/blurred1_zoom.png}
    \includegraphics[width=\textwidth]{figures/Numerical_Tests/Network_Deblurring/blurred2_zoom.png}
    \caption{}
    \end{subfigure}
    \begin{subfigure}{0.19\textwidth}
    \includegraphics[width=\textwidth]{figures/Numerical_Tests/Network_Deblurring/deblurred1_zoom.png}
    \includegraphics[width=\textwidth]{figures/Numerical_Tests/Network_Deblurring/deblurred2_zoom.png}
    \caption{}
    \end{subfigure}
    \begin{subfigure}{0.19\textwidth}
    \includegraphics[width=\textwidth]{figures/Numerical_Tests/Network_Deblurring/graphLa1_zoom.png}
    \includegraphics[width=\textwidth]{figures/Numerical_Tests/Network_Deblurring/graphLa2_zoom.png}
    \caption{}
    \end{subfigure}
    \begin{subfigure}{0.19\textwidth}
    \includegraphics[width=\textwidth]{figures/Numerical_Tests/Network_Deblurring/graphLa_residual1_zoom.png}
    \includegraphics[width=\textwidth]{figures/Numerical_Tests/Network_Deblurring/graphLa_residual2_zoom.png}
    \caption{}
    \end{subfigure}
\caption{Zoomed details for image deblurring and denoising with neural network initialization: (a) original image, (b) blurred and noisy image (PSNR = 19.97 (top) / 19.51 (bottom), SSIM = 0.3512 / 0.3023), (c) network reconstruction (PSNR = 21.91 / 19.88, SSIM = 0.4860 / 0.3952), (d) standard iterative scheme (PSNR = 22.70 / 21.53, SSIM = 0.5313 / 0.5169), (e) final reconstruction (PSNR = 22.71 / 22.09, SSIM = 0.5325 / 0.5727).}
    \label{fig:network_test_zoom}
\end{figure}

\begin{table}
\centering
   \begin{subtable}{0.6\textwidth}
   \begin{tabular}{c|c|c|c|c}
         & RRE & PSNR & SSIM & GMSD \\ \hline
        blurred  & 0.2154 & 20.12 & 0.3489 & 0.2733 \\ 
		  $\bx_{0}^{\delta}$ & 0.1719 & 22.08 & 0.4741 & 0.1878 \\
        $\bx_{1}^{\delta}$ & 0.1636 & 22.51 & 0.4993 & 0.1802 \\
		$\bx_{n_{\mathrm{std}}}^{\delta}$  & 0.1600 & 22.70 & 0.5123 & 0.1792 \\ \hline
		  $\bx_{n_{\mathrm{std}}+1}^{\delta}$ & 0.1598 & 22.71 & 0.5129 & 0.1781\\
        $\bx_{n_{\mathrm{std}}+2}^{\delta}$ & 0.1597 & 22.71 & 0.5134 & 0.1774 \\
        $\bx_{n_{\mathrm{std}}+3}^{\delta}$ & 0.1597 & 22.72 & 0.5137 & 0.1769 \\
	\end{tabular}
    \caption{}
    \end{subtable}
    \begin{subtable}{0.6\textwidth}
    \begin{tabular}{c|c|c|c|c}
         & RRE & PSNR & SSIM & GMSD \\ \hline
        blurred  & 0.1538 & 23.11 & 0.6398 & 0.2520 \\ 
		  $\bx_{0}^{\delta}$ & 0.1239 & 24.98 & 0.7369 & 0.2046 \\
        $\bx_{1}^{\delta}$ & 0.1155 & 25.59 & 0.7549 & 0.1944 \\
		$\bx_{n_{\mathrm{std}}}^{\delta}$  & 0.1111 & 25.93 & 0.7670 & 0.1812  \\ \hline
		$\bx_{n_{\mathrm{std}}+1}^{\delta}$ & 0.1101 & 26.01 & 0.7697 & 0.1752 \\
        $\bx_{n_{\mathrm{std}}+2}^{\delta}$ & 0.1093 & 26.08 & 0.7718 & 0.1704 \\
        $\bx_{n_{\mathrm{std}}+3}^{\delta}$ & 0.1086 & 26.13 & 0.7736 & 0.1667 \\
	\end{tabular}
    \caption{}
    \end{subtable}
	
	\caption{Performance metrics for deblurring and denoising with neural network initialization: RRE, PSNR, SSIM, and GMSD of the blurred image, the network reconstruction, and the iterative graph Laplacian reconstructions, where the standard iteration is used above the horizontal line and the error-equation stage of the mixed scheme is used below the line. Table (a) shows measures for the first test image (House Front) and (b) for the second test image (Windows with Flowers).}
    \label{Table:network_test}
\end{table}


\section*{Conclusions}
We have introduced three iterated graph Laplacian regularization schemes for linear image restoration problems: a standard iteration, an error-equation iteration, and a mixed strategy combining the two. In all cases, the graph Laplacian is updated throughout the reconstruction process, allowing the regularizer to progressively incorporate more accurate structural information about the unknown image.

We proved stability and convergence results in the vanishing-noise regime under a priori parameter and stopping rules. The numerical experiments on deblurring and computed tomography confirm that updating the graph improves the single-step graphLa\texttt{+}$\Psi$ reconstruction. The standard iterations mainly improve the overall reconstruction quality, whereas the error-equation stage can further enhance local details once the standard scheme has reached a plateau.


\appendix

\section{Joint coercivity assumption}\label{sec:appendix-coercivity}

In this section, we verify that the uniform joint coercivity
\Cref{eq:joint-coercivity} of \Cref{hyp:joint-coercivity} holds for two
representative model operators: a periodic deblurring model and a full-angle CT
model. The numerical experiments of \Cref{ssec:CT} instead use a
sparse-angle CT model, which is more challenging and better showcases the
reconstruction properties of the iterative methods introduced here. Uniform
joint coercivity can also be established for that operator, but the argument is
considerably more involved and lies outside the scope of this presentation; we
therefore restrict the analysis to the standard full-angle case.



By the contrast cap in \Cref{eq:gaussian-weights}, every graph input
$\bz\in X$ and every spatial edge satisfy
\[
0<\omega_*\le \omega_\bz(i,j)\le 1,
\qquad
\omega_*:=\exp\!\left(-\frac{B_{\rm c}^2}{\sigma^2}\right).
\]
Thus the estimates below are uniform over all graph inputs. When the graph input takes values in an interval
of length at most $B_{\rm c}$, the contrast cap has no effect.


\begin{remark}\label{rem:finite-dim-coercivity}
For a fixed discretization level $N$, the verification of the joint coercivity is elementary. If the spatial graph is connected, then the positive edge weights imply
\[
\ker(\Delta_\bz)=\operatorname{span}\{\bone\}
\qquad \text{for every }\bz\in X.
\]
Therefore, if $K\bone\ne0$, then $\ker K\cap\ker\Delta_\bz=\{0\}$ for every $\bz$.
Consequently there exists $C_{J,N}>0$, independent of $\bz$ but possibly
dependent on $N$, such that
\[
\norm{\bu}_2\le C_{J,N}\bigl(\norm{K\bu}_Y+\norm{\Delta_\bz\bu}_q\bigr).
\]
For each matrix operator used in the experiments (the deblurring and CT
matrices), it is straightforward to check that $K\bone\ne0$. This alone,
however, does not yield the dimension-free constant of \Cref{hyp:joint-coercivity},
since the constant $C_{J,N}$ may still diverge as $N\to\infty$.
\end{remark}

\subsection{Dimension-free estimate for a periodic blur model}\label{sec:appendix-deblurring}

We next give a genuinely joint estimate for a simple blurring model operator. Low frequencies are controlled by the blur, while the graph Laplacian controls the remaining frequencies. The result is stated for $q=2$.

Let $P_N=(\mathbb Z/n\mathbb Z)^2$, $N=n^2$, and let $K_N\colon X_N\to X_N$ be the periodic convolution operator generated by a fixed pixel-scale PSF $h\in\ell^1(\mathbb Z^2)$:
\[
K_N\bu=h*\bu.
\]
In this subsection, both the neighbor relation and the spatial distance in the
graph construction are understood periodically on $P_N$.

Let
\[
\widehat h(\theta):=\sum_{r\in\mathbb Z^2}h(r)\,e^{-ir\cdot\theta},
\qquad \theta\in[-\pi,\pi]^2.
\]
We assume only that
\[
\widehat h(0)=\sum_r h(r)\ne0.
\]
In particular, every normalized PSF satisfies this assumption. The symbol $\widehat h$ may vanish away from the origin.

Let $L_N$ be the periodic nearest-neighbor grid Laplacian,
\[
(L_N\bu)(i)=\sum_{j\sim i}(\bu(i)-\bu(j)),
\]
where $j\sim i$ denotes the four nearest neighbors. Its Fourier symbol is
\[
\lambda(\theta)=4-2\cos\theta_1-2\cos\theta_2.
\]
We assume that the graph used to build $\Delta_\bz$ contains the nearest-neighbor grid edges. Since the corresponding graph weights are bounded below by $\omega_*$, the Dirichlet forms satisfy
\begin{equation}\label{eq:graph-form-comparison}
\langle \Delta_\bz\bu,\bu\rangle
\ge \omega_*\langle L_N\bu,\bu\rangle,
\end{equation}
where the inner product is the normalized Euclidean inner product on $X_N$.

\begin{proposition}\label{prop:dim-free-deblurring}
Under the assumptions above, there exists a constant $C_J>0$, independent of $N$ and of $\bz$, such that
\[
\norm{\bu}_2\le C_J\bigl(\norm{K_N\bu}_2+\norm{\Delta_\bz\bu}_2\bigr)
\qquad\text{for every }\bu\in X_N.
\]
\end{proposition}

\begin{proof}
Since $\widehat h$ is continuous and $\widehat h(0)\ne0$, there exist $\tau>0$ and $\kappa>0$, independent of $N$, such that
\[
|\widehat h(\theta)|\ge \kappa
\qquad\text{whenever }\lambda(\theta)\le\tau.
\]
Let $P_\tau$ be the Fourier projection onto the modes with $\lambda(\theta)\le\tau$, and write
\[
\bu_\ell=P_\tau\bu,
\qquad
\bu_h=(I-P_\tau)\bu.
\]
On the low-frequency space, $K_N$ is uniformly bounded below:
\[
\norm{\bu_\ell}_2\le \kappa^{-1}\norm{K_N\bu_\ell}_2.
\]
By Young's inequality, $\norm{K_N\bv}_2\le M_K\norm{\bv}_2$ with $M_K=\norm{h}_{\ell^1}$, so
\begin{equation}\label{eq:low-bound-blur}
\norm{\bu_\ell}_2
\le \kappa^{-1}\norm{K_N\bu}_2+\kappa^{-1}M_K\norm{\bu_h}_2.
\end{equation}
For the high-frequency part, the definition of $P_\tau$ gives
\[
\tau\norm{\bu_h}_2^2\le \langle L_N\bu,\bu\rangle.
\]
Using \Cref{eq:graph-form-comparison} and Cauchy--Schwarz,
\begin{equation}\label{eq:high-bound-blur}
\norm{\bu_h}_2^2
\le (\tau\omega_*)^{-1}\langle\Delta_\bz\bu,\bu\rangle
\le (\tau\omega_*)^{-1}\norm{\Delta_\bz\bu}_2\norm{\bu}_2.
\end{equation}
Combining \Cref{eq:low-bound-blur} with $\norm{\bu}_2\le\norm{\bu_\ell}_2+\norm{\bu_h}_2$ gives
\[
\norm{\bu}_2
\le \kappa^{-1}\norm{K_N\bu}_2
+\bigl(1+\kappa^{-1}M_K\bigr)\norm{\bu_h}_2.
\]
Insert \Cref{eq:high-bound-blur}. With
\[
a:=\norm{\bu}_2,
\qquad b:=\norm{K_N\bu}_2,
\qquad c:=\norm{\Delta_\bz\bu}_2,
\]
we obtain
\[
a\le \kappa^{-1}b+
\bigl(1+\kappa^{-1}M_K\bigr)(\tau\omega_*)^{-1/2}(ca)^{1/2}.
\]
Set
\[
A:=\bigl(1+\kappa^{-1}M_K\bigr)(\tau\omega_*)^{-1/2}.
\]
Young's inequality gives $A(ca)^{1/2}\le \frac12 a+\frac12 A^2c$. Therefore
\[
a\le 2\kappa^{-1}b+A^2c.
\]
Equivalently,
\[
\norm{\bu}_2\le C_J\bigl(\norm{K_N\bu}_2+\norm{\Delta_\bz\bu}_2\bigr),
\]
with $C_J=\max\{2\kappa^{-1},A^2\}$, which is independent of $N$ and $\bz$.
\end{proof}


\subsection{Dimension-free estimate for full-angle CT}\label{sec:appendix-CT}

For CT, the cleanest dimension-free example is full-angle tomography equipped
with the natural filtered data norm. At the continuous level, the relevant
estimate is the Plancherel identity for the Radon transform. More precisely, the
Fourier slice theorem gives
\[
\mathcal F_s(\mathcal R f)(\theta,\rho)
=
\widehat f(\rho\cos\theta,\rho\sin\theta),
\]
where $\mathcal R$ is the Radon transform and $\mathcal F_s$ denotes the
one-dimensional Fourier transform in the detector variable. Hence, by
Plancherel's theorem and the polar-coordinate change of variables in frequency
space,
\[
\norm{f}_{L^2(\R^2)}^2
=
c_{\mathcal F}
\int_0^\pi\!\int_\R
|\rho|\,|\mathcal F_s(\mathcal R f)(\theta,\rho)|^2\,d\rho\,d\theta ,
\]
where $c_{\mathcal F}>0$ depends only on the Fourier-transform normalization; see,
for example, \cite{natterer2001}. Absorbing this
constant into the definition of the data norm, set
\[
\norm{g}_Y^2:=
c_{\mathcal F}
\int_0^\pi\!\int_\R
|\rho|\,|\mathcal F_s g(\theta,\rho)|^2\,d\rho\,d\theta .
\]
Then
\[
\norm{f}_{L^2(\R^2)}=\norm{\mathcal R f}_Y .
\]

At the continuous level, the above identity 
is an isometry; in particular, the coercivity bound
$\norm{f}_{L^2(\R^2)}\le\norm{\mathcal R f}_Y$ holds with constant $1$. The
finite-dimensional operator $K_N$ and the discrete data norm $\norm{\cdot}_Y$ are
obtained by discretizing $\mathcal R$ and by applying quadrature to the filtered
integral over $(\theta,\rho)$. When the discretization is consistent and the
quadrature convergent, this exact isometry passes to the discrete level as a norm
equivalence whose constants remain bounded as the mesh is refined; in particular
the lower bound
\begin{equation}\label{eq:ct-filtered-lower-bound}
\norm{\bu}_2\le C_R\norm{K_N\bu}_Y,
\end{equation}
holds with $C_R$ independent of $N$. We take \cref{eq:ct-filtered-lower-bound}
as an assumption on the discretized CT model.

\begin{proposition}\label{prop:dim-free-CT}
Assume that $K_N$ is a full-angle CT discretization satisfying
\Cref{eq:ct-filtered-lower-bound}, and that
\[
\norm{K_N\bu}_Y\le M_K\norm{\bu}_2
\]
with $M_K$ independent of $N$. Then, for every $q\in[1,2]$, the joint coercivity
condition holds on $X_N$ with a constant independent of $N$ and $\bz$:
\[
\norm{\bu}_2
\le
C_R\norm{K_N\bu}_Y
\le
C_R\bigl(\norm{K_N\bu}_Y+\norm{\Delta_\bz\bu}_q\bigr).
\]
\end{proposition}

\begin{proof}
The claim is immediate from 
\Cref{eq:ct-filtered-lower-bound}.
\end{proof}





\bibliographystyle{plain}
\bibliography{iterated_graph_Laplacian.bib}
\end{document}
